\documentclass[11pt,a4paper]{report}

% ==============================================================================
% Compatibilidad dual XeTeX (Tectonic) y pdfLaTeX (Overleaf / Lab Debian)
% ==============================================================================
\usepackage{iftex}
\ifPDFTeX
  \usepackage[utf8]{inputenc}
  \usepackage[T1]{fontenc}
  \usepackage{lmodern}
\else
  \usepackage{fontspec}
\fi

\usepackage[spanish,es-tabla,es-noquoting,es-noshorthands]{babel}
\usepackage[margin=2.4cm,headheight=15pt]{geometry}
\usepackage{microtype}
\usepackage{parskip}
\usepackage{setspace}
\onehalfspacing
\setlength{\emergencystretch}{3em}

% Colores corporativos y academicos (estilo UPM / ingenieria)
\usepackage[table,dvipsnames]{xcolor}
\definecolor{upmblue}{HTML}{003865}
\definecolor{upmaccent}{HTML}{0284C7}
\definecolor{darkslate}{HTML}{0F172A}
\definecolor{codebg}{HTML}{F8FAFC}
\definecolor{codeframe}{HTML}{CBD5E1}
\definecolor{okgreen}{HTML}{15803D}
\definecolor{warnamber}{HTML}{B45309}
\definecolor{errred}{HTML}{B91C1C}

% Graficos, tablas y diagramas vectoriales TikZ
\usepackage{graphicx}
\usepackage{booktabs}
\usepackage{tabularx}
\usepackage{longtable}
\usepackage{multirow}
\usepackage{array}
\usepackage{float}
\usepackage{caption}
\captionsetup{font=small,labelfont={bf,color=upmblue}}

\usepackage{tikz}
\usetikzlibrary{arrows.meta,positioning,shapes.geometric,fit,backgrounds,calc}

% Cajas destacadas para decisiones de ingenieria, notas y resultados
\usepackage[most]{tcolorbox}
\newtcolorbox{decisionbox}[1]{
  enhanced,
  breakable,
  colback=upmblue!4,
  colframe=upmblue,
  fonttitle=\bfseries,
  title={#1},
  left=8pt,right=8pt,top=6pt,bottom=6pt,
  boxrule=0.8pt,
  arc=3pt
}

\newtcolorbox{outcomebox}[1]{
  enhanced,
  breakable,
  colback=okgreen!5,
  colframe=okgreen!85!black,
  fonttitle=\bfseries,
  title={#1},
  left=8pt,right=8pt,top=6pt,bottom=6pt,
  boxrule=0.8pt,
  arc=3pt
}

\newtcolorbox{warnbox}[1]{
  enhanced,
  breakable,
  colback=warnamber!6,
  colframe=warnamber!90!black,
  fonttitle=\bfseries,
  title={#1},
  left=8pt,right=8pt,top=6pt,bottom=6pt,
  boxrule=0.8pt,
  arc=3pt
}

% Configuracion de bloques de codigo (HCL, YAML, XML, Bash)
\usepackage{listings}
\lstdefinestyle{cnvrcode}{
  backgroundcolor=\color{codebg},
  basicstyle=\ttfamily\footnotesize,
  breakatwhitespace=false,
  breaklines=true,
  captionpos=b,
  commentstyle=\color{gray!80!black}\itshape,
  frame=single,
  rulecolor=\color{codeframe},
  keepspaces=true,
  keywordstyle=\color{upmblue}\bfseries,
  numbers=left,
  numbersep=7pt,
  numberstyle=\tiny\color{gray},
  showspaces=false,
  showstringspaces=false,
  showtabs=false,
  stringstyle=\color{okgreen},
  tabsize=2,
  xleftmargin=12pt,
  framexleftmargin=12pt
}

\lstdefinelanguage{HCL}{
  morekeywords={terraform,required_providers,provider,variable,resource,data,output,locals,module,true,false,null,for_each,count,depends_on,validation},
  sensitive=true,
  morecomment=[l]{\#},
  morecomment=[l]{//},
  morecomment=[s]{/*}{*/},
  morestring=[b]"
}

\lstdefinelanguage{YAML}{
  morekeywords={true,false,null,apiVersion,kind,metadata,spec,name,namespace,labels,annotations,containers,initContainers,volumes,ports,selector,rules,ingress,egress,policyTypes},
  sensitive=true,
  morecomment=[l]{\#},
  morestring=[b]",
  morestring=[b]'
}

\lstset{style=cnvrcode}

% Cabeceras, titulos e hipervinculos
\usepackage{titlesec}
\titleformat{\chapter}[display]
  {\normalfont\huge\bfseries\color{upmblue}}
  {\filright\normalsize\color{upmaccent}\MakeUppercase{\chaptertitlename}\ \thechapter}{4pt}{\Huge}
\titlespacing*{\chapter}{0pt}{-10pt}{18pt}
\titleformat{\section}{\Large\bfseries\color{upmblue}}{\thesection}{0.8em}{}
\titleformat{\subsection}{\large\bfseries\color{darkslate}}{\thesubsection}{0.8em}{}

\usepackage{fancyhdr}
\pagestyle{fancy}
\fancyhf{}
\fancyhead[L]{\small\textcolor{gray}{CNVR / ARSN (2026-27) --- Trabajo Final}}
\fancyhead[R]{\small\textcolor{gray}{Grupo 11 $\cdot$ K8s + VNX + Terraform}}
\fancyfoot[C]{\small\thepage}
\renewcommand{\headrulewidth}{0.4pt}

\usepackage{enumitem}
\setlist{nosep,topsep=4pt}

\usepackage{xurl}
\usepackage[colorlinks=true,linkcolor=upmblue,urlcolor=upmaccent,citecolor=upmblue]{hyperref}

\begin{document}

% ==============================================================================
% PORTADA
% ==============================================================================
\begin{titlepage}
  \centering
  \vspace*{0.2cm}
  {\Large\bfseries\textcolor{upmblue}{UNIVERSIDAD POLITÉCNICA DE MADRID}}\\[0.2cm]
  {\large Escuela Técnica Superior de Ingenieros de Telecomunicación}\\[0.1cm]
  {\normalsize Departamento de Ingeniería de Sistemas Telemáticos (DIT)}\\[0.25cm]
  {\small\textsc{Máster Universitario en Ingeniería de Redes y Servicios Telemáticos (MUIRST)}}

  \vspace{0.8cm}
  \textcolor{upmaccent}{\rule{\linewidth}{1.5pt}}\\[0.35cm]
  {\large\bfseries\textcolor{darkslate}{Trabajo Final de Asignatura}}\\[0.25cm]
  {\huge\bfseries\textcolor{upmblue}{Despliegue Automático de una Aplicación Escalable sobre Kubernetes y VNX Orquestado con Terraform}}\\[0.3cm]
  {\normalsize Memoria Técnica, Fundamentos de Redes en Linux, Arquitectura Kubernetes y Decisiones de Ingeniería}\\[0.3cm]
  \textcolor{upmaccent}{\rule{\linewidth}{1.5pt}}

  \vspace{0.85cm}
  \begin{tabular}{rl}
    \textbf{Asignaturas:} & Computación en Nube y Virtualización de Redes y Servicios (CNVR) \\
                          & Arquitectura de Redes y Servicios en la Nube (ARSN) \\[0.2cm]
    \textbf{Autor:}       & Carlos Aparicio de Santiago \\[0.15cm]
    \textbf{Grupo:}       & Grupo 11 (\texttt{/servers-G11}) \\[0.15cm]
    \textbf{Curso:}       & 2026--2027 \\[0.15cm]
    \textbf{Fecha:}       & Noviembre de 2026
  \end{tabular}

  \vfill
  \begin{tcolorbox}[colback=codebg,colframe=codeframe,width=0.96\linewidth,arc=3pt,boxrule=0.6pt]
    \footnotesize
    \textbf{Resumen Ejecutivo:} Este documento detalla el diseño, la fundamentación teórica y la automatización extremo a extremo de una infraestructura híbrida compuesta por un clúster de Kubernetes (con CNI primario \textbf{Calico} y meta-CNI \textbf{Multus}) interconectado a nivel de enlace mediante una VLAN IEEE 802.1Q compartida (\texttt{VlanNet}) con contenedores \textbf{LXC} gestionados por \textbf{VNX} (servidor de base de datos \textbf{MariaDB} y servidor centralizado de \textbf{Rsyslog}). Todo el ciclo de vida de los recursos de red, secretos SSH, plantillas HTML dinámicas, despliegues escalables, balanceo Ingress y políticas de seguridad \textit{Zero-Trust} (\texttt{NetworkPolicy}) se orquesta de manera declarativa mediante \textbf{Terraform}.
  \end{tcolorbox}
\end{titlepage}

\pagenumbering{roman}
\tableofcontents
\clearpage
\pagenumbering{arabic}

% ==============================================================================
% CAPITULOS MODULARES
% ==============================================================================
\chapter{Introducción y Arquitectura Global del Sistema}

\section{Contexto y Objetivos del Trabajo Final}

La evolución de los servicios telemáticos modernos hacia arquitecturas nativas en la nube (\textit{cloud-native}) exige combinar tres pilares tecnológicos fundamentales: la contenedorización orquestada para lograr elasticidad y auto-reparación (\textbf{Kubernetes}), la virtualización avanzada de redes en sistemas Linux para aislar y acoplar planos de tráfico heterogéneos (\textbf{Calico}, \textbf{Multus}, \textbf{Open vSwitch} y \textbf{VLANs 802.1Q}), y la Infraestructura como Código (\textit{Infrastructure as Code}, \textbf{Terraform}) para garantizar despliegues reproducibles, declarativos y auditables.

Partiendo de la base establecida en las prácticas previas de la asignatura ---la orquestación declarativa sobre OpenStack y Kubernetes con Calico en la \textbf{Práctica 1.4}, la virtualización de redes con Open vSwitch y VLANs en la \textbf{Práctica 2.1}, y el análisis de comunicaciones de Pods, volúmenes, servicios, MetalLB, Ingress y Multus en la \textbf{Práctica 3.1}---, este trabajo final aborda la construcción y automatización completa de una plataforma híbrida en la que conviven aplicaciones contenerizadas en Kubernetes con servicios de infraestructura persistentes desplegados sobre contenedores ligeros de sistema (\textbf{LXC}) mediante \textbf{VNX}.

Los objetivos específicos alcanzados en este proyecto son:
\begin{enumerate}
  \item \textbf{Automatización integral en un único flujo (\texttt{01\_deploy\_all.sh})}: desde la generación criptográfica de pares de claves SSH, el arranque del clúster de Kubernetes con Calico (\path{k8s_lab-calico-v02}) y la provisión del nodo dedicado \texttt{terraform} (\path{k8s_lab-terraform.xml} con carpeta compartida \texttt{/shared}), hasta la creación de subinterfaces VLAN en los nodos \textit{worker}, el despliegue automático de contenedores LXC en VNX y la convergencia declarativa con Terraform.
  \item \textbf{Interconexión híbrida en Capa 2 mediante Multus y VLAN 802.1Q}: integración de los Pods de Kubernetes y los contenedores LXC externos sobre una misma red virtual troncal (\texttt{VlanNet}, VLAN \texttt{1000}, subred \texttt{10.1.2.0/24}), resolviendo la asignación dinámica de direcciones IP para permitir el escalado horizontal sin colisiones.
  \item \textbf{Gestión segura sin contraseñas mediante Kubernetes \texttt{Secrets}}: separación estricta de credenciales en dos pares de claves SSH independientes (acceso exterior hacia administración y administración hacia servidores internos), gestionados íntegramente desde HCL en Terraform sin recurrir a comandos imperativos de \texttt{kubectl}.
  \item \textbf{Personalización dinámica del plano de aplicación}: renderizado automático de páginas HTML individualizadas con el nombre e identidad de cada Pod mediante la combinación de un \texttt{ConfigMap} de solo lectura, la \textit{Downward API} de Kubernetes, un volumen compartido \texttt{emptyDir} y un \texttt{initContainer}.
  \item \textbf{Exposición controlada y seguridad \textit{Zero-Trust}}: publicación HTTP mediante un controlador \textbf{Ingress NGINX} en la ruta de grupo \texttt{/servers-G11} hacia un \texttt{Service ClusterIP}, acceso SSH administrativo vía \texttt{Service NodePort} conservando la IP de origen (\texttt{externalTrafficPolicy: Local}) y aislamiento estricto de entrada y salida mediante tres políticas \texttt{NetworkPolicy} aplicadas por Calico.
\end{enumerate}

\section{Topología Física y Lógica del Escenario}

La Figura~\ref{fig:topologia-global} ilustra la arquitectura completa implementada siguiendo el escenario de laboratorio de la Práctica 1.4 ampliado con la red VLAN compartida hacia VNX. El sistema se articula sobre los nodos de \texttt{k8s\_LAN} (incluyendo el nodo proveedor \texttt{terraform} en \texttt{10.10.10.13}) y el escenario complementario de contenedores LXC:

\begin{figure}[H]
\centering
\resizebox{\textwidth}{!}{%
\begin{tikzpicture}[
  font=\sffamily\scriptsize,
  box/.style={draw=upmblue, rounded corners=3pt, fill=white, thick, align=center, inner sep=5pt},
  tfbox/.style={draw=Purple!80!black, rounded corners=3pt, fill=Purple!7, thick, align=center, inner sep=4.5pt},
  pod/.style={draw=upmaccent, rounded corners=6pt, fill=upmaccent!8, thick, align=center, minimum width=2.4cm, minimum height=0.85cm},
  svc/.style={draw=okgreen!80!black, rounded corners=8pt, fill=okgreen!8, thick, align=center, inner sep=4pt},
  lxc/.style={draw=warnamber!90!black, cylinder, shape border rotate=90, aspect=0.22, fill=warnamber!10, thick, align=center, minimum width=2.3cm, minimum height=1.3cm},
  line/.style={-{Stealth[length=2mm]}, thick, color=darkslate},
  tfline/.style={-{Stealth[length=2mm]}, thick, dashed, color=Purple!80!black},
  vlanline/.style={very thick, color=upmaccent}
]

  % Host Exterior, Router r1 y Nodo Terraform
  \node[box, fill=gray!10, minimum width=2.5cm, minimum height=1.2cm] (host) at (-8.1, 0.8)
    {\textbf{Host Laboratorio}\\Red \texttt{user\_LAN}\\\texttt{10.10.20.2/24}};

  \node[box, fill=upmblue!8, minimum width=2.0cm, minimum height=1.0cm] (r1) at (-5.1, 0.8)
    {\textbf{Router \texttt{r1}}\\\texttt{.20.1 $\leftrightarrow$ .10.1}};

  \node[tfbox, minimum width=3.5cm, minimum height=1.15cm] (tfvm) at (-6.5, -2.1)
    {\textbf{Nodo VNX \texttt{terraform}}\\\texttt{10.10.10.13} (\texttt{/shared})};

  % Caja Cluster Kubernetes
  \node[draw=upmblue, dashed, thick, rounded corners=5pt, fill=upmblue!2, minimum width=8.6cm, minimum height=7.0cm,
        label={[anchor=north west, yshift=-1pt, font=\sffamily\bfseries\scriptsize, text=upmblue]north west:Clúster Kubernetes (\texttt{.10.10--.12}, Calico + Multus)}] (k8sbox) at (1.5, 0.25) {};

  % Ingress y Services bien separados
  \node[svc] (ing) at (-1.4, 1.6) {\textbf{Ingress NGINX}\\\texttt{10.10.10.20}\\\texttt{/servers-G11}};
  \node[svc] (svcweb) at (1.3, 1.6) {\textbf{web-service}\\ClusterIP \texttt{:8080}};
  \node[svc] (svcadm) at (0.2, -1.8) {\textbf{admin-service}\\NodePort \texttt{:30022} (\texttt{Local})};

  % Pods Web y Admin
  \node[pod] (s1) at (4.1, 2.4) {\textbf{Pod Web 1}\\\texttt{nginx-ssh} (\texttt{:80,:22})};
  \node[pod] (s2) at (4.1, 1.2) {\textbf{Pod Web 2}\\\texttt{nginx-ssh} (\texttt{:80,:22})};
  \node[pod] (s3) at (4.1, 0.0) {\textbf{Pod Web 3}\\\texttt{nginx-ssh} (\texttt{:80,:22})};
  \node[pod, fill=upmblue!12, draw=upmblue] (adm) at (4.1, -1.8) {\textbf{Pod \texttt{admin}}\\\texttt{k8s-worker1} (\texttt{10.1.2.10})};

  % Bus VLAN 1000 entre ambas cajas
  \draw[vlanline] (6.4, 3.1) -- (6.4, -3.25) node[below=2pt, font=\sffamily\bfseries\scriptsize, text=upmaccent] {VLAN 1000 (\texttt{VlanNet})};

  % Caja VNX
  \node[draw=warnamber!90!black, dashed, thick, rounded corners=5pt, fill=warnamber!3, minimum width=3.1cm, minimum height=7.0cm,
        label={[anchor=north, font=\sffamily\bfseries\scriptsize, text=warnamber!90!black]north:VNX (LXC)}] (vnxbox) at (8.7, 0.25) {};

  \node[lxc] (bbdd) at (8.7, 1.5) {\textbf{LXC \texttt{bbdd}}\\MariaDB \texttt{:3306}\\\texttt{10.1.2.100}};
  \node[lxc] (logs) at (8.7, -1.3) {\textbf{LXC \texttt{logs}}\\Rsyslog \texttt{:514} + DHCP\\\texttt{10.1.2.101}};

  % Conexiones HTTP, SSH exterior y Terraform
  \draw[line] (host) -- (r1);
  \draw[tfline] (host.south) |- node[pos=0.25, right, font=\sffamily\tiny] {\texttt{/shared}} (tfvm.west);
  \draw[tfline] (tfvm.east) -- node[below, font=\sffamily\tiny] {API K8s \texttt{:6443}} (k8sbox.west |- tfvm.east);

  \draw[line] (r1.north east) |- node[pos=0.58, above, font=\sffamily\tiny] {HTTP \texttt{:80}} (ing.west);
  \draw[line] (r1.south east) |- node[pos=0.56, above, font=\sffamily\tiny] {SSH (\texttt{.10.11:30022})} (svcadm.west);

  \draw[line] (ing) -- (svcweb);
  \draw[line] (svcweb.east) -- (s1.west);
  \draw[line] (svcweb.east) -- (s2.west);
  \draw[line] (svcweb.east) -- (s3.west);
  \draw[line] (svcadm) -- (adm);

  % Conexiones a la VLAN 1000
  \draw[vlanline] (s1.east) -- (6.4, 2.4);
  \draw[vlanline] (s2.east) -- (6.4, 1.2);
  \draw[vlanline] (s3.east) -- (6.4, 0.0);
  \draw[vlanline] (adm.east) -- (6.4, -1.8);
  \draw[vlanline] (6.4, 1.5) -- (bbdd.west);
  \draw[vlanline] (6.4, -1.3) -- (logs.west);

\end{tikzpicture}%
}
\caption{Topología híbrida desplegada: nodo proveedor \texttt{terraform} (\texttt{.13}), clúster Kubernetes (Calico + Multus) e infraestructura VNX (LXC) interconectados mediante la VLAN 1000 sobre \texttt{VlanNet}.}
\label{fig:topologia-global}
\end{figure}

\subsection{Planos de Red y Direccionamiento Completo}

Para comprender el flujo extremo a extremo de cada paquete sin ambigüedades, es esencial distinguir los seis rangos IP que operan simultáneamente sobre la infraestructura virtualizada (Tabla~\ref{tab:planos-red}):

\begin{table}[H]
\centering
\small
\renewcommand{\arraystretch}{1.25}
\begin{tabularx}{\textwidth}{l l l X}
\toprule
\textbf{Red / Plano} & \textbf{Prefijo CIDR} & \textbf{Interfaz} & \textbf{Propósito y Elementos Conectados} \\
\midrule
\texttt{user\_LAN} & \texttt{10.10.20.0/24} & Host $\leftrightarrow$ \texttt{r1} & Red de acceso del usuario en el laboratorio. El host (\texttt{10.10.20.2}) accede a través del router \texttt{r1} (\texttt{10.10.20.1}). \\
\texttt{k8s\_LAN} & \texttt{10.10.10.0/24} & \texttt{e1} (nodos VNX) & Red troncal: \texttt{r1} (\texttt{.1}), \texttt{k8s-master} (\texttt{.10}), \texttt{k8s-worker1} (\texttt{.11}), \texttt{k8s-worker2} (\texttt{.12}), nodo proveedor \textbf{\texttt{terraform} (\texttt{.13})} y pool MetalLB (\texttt{.20--.30}). \\
Red Pods Calico & \texttt{10.233.64.0/18} (o \texttt{/16}) & \texttt{eth0} $\leftrightarrow$ \texttt{cali*} & Red primaria CNI de Kubernetes gestionada por Calico (túneles IP-in-IP inter-nodo); gobernada por \texttt{NetworkPolicy}. \\
Red Services K8s & \texttt{10.96.0.0/12} & Virtual (\texttt{iptables}) & Direcciones \texttt{ClusterIP} virtuales programadas por \texttt{kube-proxy} (no residen en ninguna interfaz física ni virtual). \\
\texttt{VlanNet} (VLAN 1000) & \texttt{10.1.2.0/24} & \texttt{net1} / \texttt{e2.1000} / \texttt{eth1} & Plano compartido de datos y gestión L2: Pod \texttt{admin} (\texttt{.10}), Pods \texttt{web} (\texttt{.50--.95} DHCP), LXC \texttt{bbdd} (\texttt{.100}) y LXC \texttt{logs} (\texttt{.101}). \\
\texttt{virbr0} (NAT) & \texttt{192.168.122.0/24} & \texttt{e9} / \texttt{eth9} & Bridge libvirt de salida a Internet para \texttt{k8s-*}, \texttt{terraform} (descarga de proveedores HCL) y contenedores LXC. \\
\bottomrule
\end{tabularx}
\caption{Resumen de los planos de red físicos, virtuales y lógicos del escenario.}
\label{tab:planos-red}
\end{table}

\section{Estructura Modular del Proyecto Entregado}

Siguiendo las mejores prácticas de ingeniería de software e Infraestructura como Código, el repositorio desarrollado en \texttt{cnvr\_final/} separa limpiamente la automatización imperativa de plataforma (\texttt{scripts/} y \texttt{vnx/}) de la orquestación declarativa de servicios (\texttt{terraform/}):

\begin{itemize}
  \item \textbf{\texttt{scripts/01\_deploy\_all.sh}}: script maestro directo e idempotente en 6 pasos (generación de los dos pares RSA-4096, arranque de \path{k8s_lab-calico-v02} junto al nodo VNX \texttt{terraform} mediante \path{k8s_lab-terraform.xml}, creación de subinterfaces \texttt{e2.1000} y demonio DHCP CNI en los \textit{workers}, despliegue del escenario complementario LXC en VNX, sincronización en \texttt{/shared} + ejecución de \texttt{terraform init/validate/apply} e instalación de \texttt{mariadb-client} sobre la VLAN).
  \item \textbf{\texttt{scripts/02\_verify\_all.sh}} y \textbf{\texttt{scripts/03\_destroy\_all.sh}}: suite automatizada de auditoría funcional/seguridad (incluyendo \texttt{terraform state list} y \texttt{terraform output}) y desmontaje ordenado del escenario.
  \item \textbf{\texttt{vnx/k8s\_lxc\_vlan.xml}} y \textbf{\texttt{vnx/conf/}}: definición VNX 2.0 de los contenedores LXC \texttt{bbdd} y \texttt{logs} conectada a \texttt{VlanNet} (inspirada en la plantilla base \path{k8s_lab-vms-vlan.xml} del laboratorio) junto con los ficheros inyectados mediante \texttt{<filetree>} y \texttt{<exec>}.
  \item \textbf{\texttt{terraform/*.tf}} y \textbf{\texttt{terraform.tfvars}}: módulos HCL (\path{k8s-provider.tf}, \path{variables.tf}, \path{terraform.tfvars}, \path{locals.tf}, \path{secrets.tf}, \path{multus.tf}, \path{configmap.tf}, \path{web.tf}, \path{admin.tf}, \path{services_ingress.tf}, \path{network_policies.tf} y \path{outputs.tf}) validados contra la versión exacta \texttt{2.23.0} del proveedor oficial \path{hashicorp/kubernetes} empleada en la Práctica 1.4.
\end{itemize}

\chapter{Fundamentos de Redes en Linux Aplicados al Escenario}

Kubernetes y VNX no implementan una pila de red propietaria: delegan toda la conmutación, el aislamiento, el etiquetado VLAN, el enrutamiento y el filtrado de paquetes en primitivas nativas del kernel de Linux. Entender qué ocurre exactamente en el kernel permite razonar con rigor sobre cada decisión tomada en la práctica.

\section{Network Namespaces y el Contenedor \texttt{pause}}

Un \textit{namespace} de Linux proporciona a un conjunto de procesos una vista aislada de un subsistema del kernel (\texttt{pid}, \texttt{mount}, \texttt{uts}, \texttt{ipc}, \texttt{cgroup}, \texttt{user} y \texttt{net}). En particular, cada \textbf{network namespace (\texttt{netns})} posee de forma completamente independiente:
\begin{itemize}
  \item Su propio juego de interfaces de red (incluido su propio \textit{loopback} \texttt{lo}).
  \item Su propia tabla de encaminamiento IPv4/IPv6 y su propia caché ARP (\textit{neighbour table}).
  \item Sus propias tablas de \texttt{netfilter}/\texttt{iptables} y estado de seguimiento de conexiones (\texttt{conntrack}).
  \item Su propio espacio de puertos TCP/UDP (\textit{sockets}).
\end{itemize}

En un nodo \textit{worker} de Kubernetes coexisten el \textbf{namespace raíz del nodo} (donde residen las interfaces de la máquina virtual \texttt{e1}, \texttt{e2}, \texttt{e9}, las subinterfaces \texttt{e2.1000} y las reglas \texttt{iptables} de \texttt{kube-proxy}/Calico) y \textbf{un network namespace independiente por cada Pod}.

Cuando el \texttt{kubelet} instancia cualquiera de nuestros Pods (por ejemplo, una réplica de \texttt{web-deployment}), arranca en primer lugar un contenedor ligero de infraestructura llamado \textbf{\texttt{pause}} (\texttt{registry.k8s.io/pause}). Su única misión es crear los namespaces \texttt{net}, \texttt{ipc} y \texttt{uts} del Pod y bloquearse indefinidamente mediante la llamada al sistema \texttt{pause()}. Todos los demás contenedores del mismo Pod ---tanto nuestros \texttt{initContainers} (\texttt{generate-html} y \texttt{setup-ssh-auth}) como el contenedor principal \texttt{nginx-ssh}--- se unen al \texttt{/proc/<PID\_pause>/ns/net} ya existente. De ahí se derivan tres consecuencias arquitectónicas directas en nuestro trabajo:
\begin{enumerate}
  \item \textbf{Identidad de red única por Pod}: las dos interfaces que configura CNI en el Pod (\texttt{eth0} por Calico y \texttt{net1} por Multus) pertenecen al \texttt{netns} del \texttt{pause} y son visibles con idéntica IP y MAC por todos los contenedores del Pod.
  \item \textbf{Persistencia de red ante reinicios}: si el proceso \texttt{nginx} o \texttt{sshd} fallase y \texttt{containerd} recrease el contenedor, el \texttt{pause} permanece vivo conservando intactas sus IPs en Calico y en la VLAN \texttt{1000}.
  \item \textbf{Coexistencia de servicios en puertos distintos}: dentro de \texttt{dgsgiros/nginx-ssh:latest}, \texttt{sshd} escucha en \texttt{0.0.0.0:22} y \texttt{nginx} en \texttt{0.0.0.0:80} dentro del mismo espacio de sockets sin conflicto.
\end{enumerate}

\section{Pares \texttt{veth}, Bridges Linux y Open vSwitch (\texttt{VlanNet})}

\subsection{Cables Virtuales \texttt{veth} e Índices Cruzados}

Dado que una interfaz de red solo puede existir en un único \texttt{netns} simultáneamente, para comunicar la red primaria de un Pod con el nodo se emplea un \textbf{par \texttt{veth} (\textit{Virtual Ethernet})}: un túnel bidireccional en memoria donde todo frame Ethernet inyectado por un extremo emerge instantáneamente por el opuesto:
\begin{itemize}
  \item En la Práctica 3.1 (con \textbf{Flannel}), el extremo dentro del Pod (\texttt{eth0@if12}) se conectaba a un extremo \texttt{vethXXXXXXXX@if2} esclavizado a un \textbf{Linux bridge} común (\texttt{cni0}), que actuaba como conmutador L2 local y puerta de enlace (\texttt{10.244.X.1/24}).
  \item En el Trabajo Final (con \textbf{Calico}), desaparece el bridge compartido \texttt{cni0}: el extremo del nodo se denomina \textbf{\texttt{caliXXXXXXXX@if2}} y NO se conecta a ningún bridge L2, sino que opera como un enlace punto a punto enrutado en Capa 3 directamente por el kernel del nodo \textit{worker}. Esto evita tormentas de difusión ARP entre Pods y permite a Calico aplicar listas de control de acceso individuales sobre cada interfaz \texttt{cali*}.
\end{itemize}

\subsection{Open vSwitch (\texttt{VlanNet}) y Etiquetado IEEE 802.1Q}

Mientras que \texttt{virbr0} es un \textit{Linux bridge} clásico empleado para proporcionar acceso NAT a Internet, la red compartida \textbf{\texttt{VlanNet}} creada por VNX en el host físico del laboratorio es un conmutador virtual \textbf{Open vSwitch (OVS)}.

El estándar \textbf{IEEE 802.1Q} inserta una etiqueta de 4 bytes entre la dirección MAC origen y el campo \textit{EtherType} de la trama Ethernet original:
\begin{itemize}
  \item \textbf{TPID (16 bits)}: valor \texttt{0x8100} que identifica una trama etiquetada 802.1Q.
  \item \textbf{TCI (16 bits)}: compuesto por 3 bits de prioridad \textit{PCP}, 1 bit \textit{DEI} y \textbf{12 bits de VLAN ID (VID)}, permitiendo hasta 4094 redes virtuales aisladas (en nuestro diseño, \textbf{VID = 1000}).
\end{itemize}

En nuestro escenario convergen dos mecanismos complementarios de etiquetado 802.1Q sobre \texttt{VlanNet}:
\begin{enumerate}
  \item \textbf{En los nodos \textit{worker} de Kubernetes (puerto \textit{trunk} + subinterfaz del kernel)}: la interfaz \texttt{e2} de \texttt{k8s-worker1} y \texttt{k8s-worker2} se conecta a \texttt{VlanNet} sin filtro de VLAN en OVS (modo \textit{trunk}). Dentro del sistema operativo de cada \textit{worker}, nuestro script crea la subinterfaz \texttt{e2.1000} mediante \texttt{iproute2}:
\begin{lstlisting}[language=bash]
ip link set e2 up
ip link add link e2 name e2.1000 type vlan id 1000
ip link set e2.1000 up
\end{lstlisting}
  Toda trama transmitida por \texttt{e2.1000} recibe en el kernel del \textit{worker} la cabecera 802.1Q con \texttt{ID=1000} antes de salir por \texttt{e2} hacia \texttt{VlanNet}; recíprocamente, al recibir por \texttt{e2} una trama con \texttt{ID=1000}, el kernel retira la etiqueta y la entrega limpia a \texttt{e2.1000}.
  \item \textbf{En los contenedores LXC de VNX (puerto \textit{access} en Open vSwitch)}: en \texttt{vnx/k8s\_lxc\_vlan.xml}, la interfaz \texttt{eth1} de los contenedores \texttt{bbdd} y \texttt{logs} incluye \texttt{<vlan><tag id="1000"/></vlan>}. VNX traduce esto en un puerto de acceso OVS (\texttt{ovs-vsctl set port ... tag=1000}): el contenedor LXC transmite tramas Ethernet normales sin etiquetar por \texttt{eth1}, y es el propio switch \texttt{VlanNet} en el host quien inserta la etiqueta \texttt{1000} al entrar al switch y la retira al entregarla al contenedor LXC.
\end{enumerate}

\section{El Driver \texttt{macvlan} y su Interacción con la Pila IP del Nodo}

Para conectar cada Pod de Kubernetes a la subinterfaz \texttt{e2.1000} sin crear decenas de pares \texttt{veth} ni puentes adicionales, Multus utiliza el plugin CNI \textbf{\texttt{macvlan}} en modo \textbf{\texttt{bridge}}.

\texttt{macvlan} permite crear múltiples interfaces virtuales hijas (\texttt{net1} dentro del \texttt{netns} de cada Pod), cada una dotada de su \textbf{propia dirección MAC única}, vinculadas directamente a una interfaz padre del nodo (\texttt{master: "e2.1000"}). Entre los cuatro modos de operación de \texttt{macvlan} (\textit{private}, \textit{vepa}, \textit{bridge} y \textit{passthru}), el modo \textbf{\texttt{bridge}} resulta imprescindible en nuestro diseño porque incorpora un puente ligero en memoria entre todas las subinterfaces hijas colgadas del mismo padre \texttt{e2.1000}. Gracias a ello, cuando el Pod \texttt{admin} (alojado en \texttt{k8s-worker1}) se comunica por la VLAN con una réplica web que el planificador de Kubernetes haya colocado también en \texttt{k8s-worker1}, la trama se conmuta directamente en memoria dentro de \texttt{e2.1000} sin necesidad de salir al exterior y rebotar (\textit{hairpinning}).

La Figura~\ref{fig:arquitectura-worker1} detalla la anatomía exacta de interfaces y espacios de red en el interior de \texttt{k8s-worker1}:

\begin{figure}[H]
\centering
\resizebox{\textwidth}{!}{%
\begin{tikzpicture}[
  font=\sffamily\scriptsize,
  ns/.style={draw=upmaccent, thick, rounded corners=4pt, fill=upmaccent!5, inner sep=6pt},
  iface/.style={draw=darkslate, fill=white, rounded corners=2pt, inner sep=3pt, align=center},
  kbox/.style={draw=errred!80!black, fill=errred!6, rounded corners=3pt, inner sep=4pt, align=center},
  ovsbox/.style={draw=warnamber!90!black, fill=warnamber!10, rounded corners=3pt, thick, inner sep=5pt, align=center},
  p2p/.style={<->, thick, color=upmblue},
  l2/.style={<->, very thick, color=upmaccent}
]

  % Marco del nodo k8s-worker1
  \node[draw=upmblue, thick, rounded corners=6pt, fill=gray!4, minimum width=14.6cm, minimum height=6.6cm,
        label={[anchor=north west, font=\sffamily\bfseries\footnotesize, text=upmblue]north west:Nodo \texttt{k8s-worker1} (netns raíz)}] (w1) at (0,0) {};

  % Netns Pod Admin
  \node[ns, minimum width=5.4cm, minimum height=2.1cm,
        label={[anchor=north, font=\sffamily\bfseries\scriptsize, text=upmblue]north:\texttt{netns} Pod \texttt{admin} (\texttt{pause})}] (podadm) at (-3.7, 1.7) {};
  \node[iface] (admeth0) at (-5.1, 1.5) {\texttt{eth0}\\IP Calico};
  \node[iface, draw=upmaccent, thick] (admnet1) at (-2.3, 1.5) {\texttt{net1} (macvlan)\\\texttt{10.1.2.10/24}};

  % Netns Pod Web
  \node[ns, minimum width=5.4cm, minimum height=2.1cm,
        label={[anchor=north, font=\sffamily\bfseries\scriptsize, text=upmblue]north:\texttt{netns} Pod \texttt{web-1} (\texttt{pause})}] (podweb) at (3.7, 1.7) {};
  \node[iface] (webeth0) at (2.2, 1.5) {\texttt{eth0}\\IP Calico};
  \node[iface, draw=upmaccent, thick] (webnet1) at (5.1, 1.5) {\texttt{net1} (macvlan)\\\texttt{10.1.2.5x/24}};

  % Interfaces cali* y motor Calico (separado limpiamente de e2.1000)
  \node[iface] (cali1) at (-5.1, -0.35) {\texttt{cali8a1f...}};
  \node[iface] (cali2) at (2.2, -0.35) {\texttt{calif39b...}};
  \node[kbox, minimum width=8.6cm] (routing) at (-2.3, -1.65)
    {\textbf{Pila IP L3 del Kernel + \texttt{iptables}}\\(\texttt{kube-proxy} DNAT + Calico Felix \texttt{NetworkPolicy})};
  \node[iface] (e1) at (-2.3, -2.85) {\texttt{e1}: \texttt{10.10.10.11/24} (\texttt{k8s\_LAN})};

  % Subinterfaz VLAN y e2 a la derecha sin solape
  \node[iface, draw=upmaccent, very thick, minimum width=3.9cm] (e21000) at (4.7, -1.55)
    {\texttt{e2.1000} (802.1Q VID 1000)\\Puente interno \texttt{macvlan bridge}};
  \node[iface, draw=upmaccent, thick] (e2) at (4.7, -2.85) {\texttt{e2} (Trunk sin IP)};

  % Switch externo VlanNet
  \node[ovsbox, minimum width=5.0cm] (vlannet) at (4.7, -4.3)
    {\textbf{Open vSwitch \texttt{VlanNet} (Host)}\\$\leftrightarrow$ \texttt{k8s-worker2} y LXC \texttt{bbdd}/\texttt{logs}};

  % Conexiones
  \draw[p2p] (admeth0) -- node[left, font=\sffamily\tiny] {veth} (cali1);
  \draw[p2p] (webeth0) -- node[left, font=\sffamily\tiny] {veth} (cali2);
  \draw[p2p] (cali1) -- (routing.north -| cali1);
  \draw[p2p] (cali2.south) |- ([yshift=3pt]routing.north east);
  \draw[p2p] (routing) -- (e1);

  \draw[l2] (admnet1.south) -- (-2.3, 0.25) -| ([xshift=-8mm]e21000.north);
  \draw[l2] (webnet1.south) -- (e21000.north -| webnet1);
  \draw[l2] (e21000) -- (e2);
  \draw[l2] (e2) -- (vlannet);

\end{tikzpicture}%
}
\caption{Separación física y lógica de planos dentro de \texttt{k8s-worker1}: el tráfico de \texttt{eth0} atraviesa las interfaces \texttt{cali*} y el cortafuegos L3 de Calico, mientras que \texttt{net1} (\texttt{macvlan}) desciende directamente en Capa 2 por \texttt{e2.1000} hacia \texttt{VlanNet}.}
\label{fig:arquitectura-worker1}
\end{figure}

\begin{warnbox}{Propiedad Crítica del Plano \texttt{macvlan} en Linux}
El tráfico emitido por un Pod a través de su interfaz secundaria \texttt{net1} (\texttt{macvlan}) \textbf{no atraviesa la pila de encaminamiento IP ni los ganchos \texttt{iptables} \texttt{FORWARD} de Calico asociados a \texttt{cali*}}. Por tanto:
\begin{itemize}
  \item Las \texttt{NetworkPolicies} de Kubernetes aíslan de forma hermética el interfaz primario \texttt{eth0} (acceso HTTP público, acceso SSH exterior desde \texttt{user\_LAN} y bloqueo de salida hacia Internet).
  \item Simultáneamente, la VLAN \texttt{1000} sobre \texttt{net1} opera como una red de gestión y datos fuera de banda (\textit{out-of-band L2 fabric}) dedicada exclusivamente a los componentes autorizados conectados a \texttt{VlanNet}.
\end{itemize}
\end{warnbox}

\section{\texttt{netfilter}, \texttt{iptables} y el Problema del SNAT en \texttt{NodePort}}

En Kubernetes, las direcciones \texttt{ClusterIP} (como la de nuestro \texttt{web-service}) no están asignadas a ninguna interfaz: existen únicamente como reglas de traducción de destino (\textbf{DNAT}) en la tabla \texttt{nat} de \texttt{iptables} mantenida por \texttt{kube-proxy}. Cuando un paquete va dirigido a \texttt{ClusterIP:8080}, la cadena \texttt{KUBE-SERVICES} salta a \texttt{KUBE-SVC-XXX}, la cual distribuye las nuevas conexiones mediante el módulo \texttt{statistic --mode random --probability} entre las cadenas \texttt{KUBE-SEP-*} (\textit{Service Endpoints}), reescribiendo la IP y puerto destino por el \texttt{PodIP:80} elegido y registrando la tupla en \texttt{conntrack}.

Un fenómeno crítico ocurre con los servicios de tipo \textbf{\texttt{NodePort}} (como \texttt{admin-service}):
\begin{enumerate}
  \item \textbf{Comportamiento por defecto (\texttt{externalTrafficPolicy: Cluster})}: Kubernetes abre el puerto \texttt{NodePort} en \emph{todos} los nodos del clúster. Si el cliente externo (\texttt{10.10.20.2}) se conecta a un nodo que no aloja el Pod (o incluso al nodo que sí lo aloja bajo política \texttt{Cluster}), \texttt{kube-proxy} marca el paquete con \texttt{KUBE-MARK-MASQ} y le aplica \textbf{SNAT (\texttt{MASQUERADE})} en \texttt{POSTROUTING}, reemplazando la IP origen real del cliente (\texttt{10.10.20.2}) por la IP interna del nodo (\texttt{10.10.10.11}). Esto garantiza que el paquete de respuesta regrese al mismo nodo para deshacer el DNAT, pero \textbf{destruye la visibilidad de la IP origen} ante cualquier \texttt{NetworkPolicy} basada en \texttt{ipBlock}.
  \item \textbf{Comportamiento configurado (\texttt{externalTrafficPolicy: Local})}: \texttt{kube-proxy} elimina el salto de \texttt{MASQUERADE} y encamina únicamente hacia \textit{endpoints} locales en ese mismo nodo (descartando el tráfico en nodos donde el Pod no esté presente). Combinado con \texttt{nodeSelector: k8s-worker1}, el paquete TCP SYN enviado desde el host (\texttt{10.10.20.2}) a \texttt{10.10.10.11:30022} sufre únicamente DNAT hacia \texttt{PodAdminIP:22}, conservando intacta su cabecera \texttt{SRC=10.10.20.2} cuando Calico evalúa la regla de cortafuegos en \texttt{cali*}.
\end{enumerate}

\chapter{Arquitectura y Recursos de Kubernetes en Profundidad}

\section{Plano de Control, Nodos Worker y Evolución de CNI: Flannel vs Calico}

El clúster de Kubernetes desplegado consta de un nodo de control (\texttt{k8s-master}, \texttt{10.10.10.10}) y dos nodos de ejecución (\texttt{k8s-worker1}, \texttt{10.10.10.11} y \texttt{k8s-worker2}, \texttt{10.10.10.12}), acompañados en \texttt{k8s\_LAN} por el nodo orquestador \texttt{terraform} (\texttt{10.10.10.13}):
\begin{itemize}
  \item \textbf{Plano de Control (\texttt{k8s-master})}: aloja \texttt{etcd} (almacén clave-valor consistente mediante consenso Raft donde persiste todo el estado del clúster, incluidos los \texttt{Secrets} y \texttt{NetworkPolicies}), \texttt{kube-apiserver} (único punto de entrada REST/JSON sobre TLS \texttt{:6443} con el que dialoga el proveedor de Terraform), \texttt{kube-scheduler} (encargado de ubicar los Pods evaluando restricciones como \texttt{nodeSelector}) y \texttt{kube-controller-manager} (ejecutor de los bucles de control de \texttt{Deployment}, \texttt{ReplicaSet} y \texttt{Endpoints}).
  \item \textbf{Plano de Datos (\texttt{k8s-worker1/2})}: ejecuta el agente \texttt{kubelet}, el motor de contenedores \texttt{containerd} (estándar CRI), \texttt{kube-proxy} y los demonios de red CNI (\texttt{calico-node} y \texttt{multus}).
\end{itemize}

\subsection{Por qué Calico Sustituye a Flannel y Cómo Opera su Túnel IP-in-IP}

En la Práctica 3.1 trabajamos sobre \textbf{Flannel}, mientras que tanto la Práctica 1.4 como el Trabajo Final emplean el escenario \path{k8s_lab-calico-v02} basado en \textbf{Calico}. La Tabla~\ref{tab:flannel-vs-calico} sintetiza las diferencias técnicas que motivan este cambio:

\begin{table}[H]
\centering
\small
\renewcommand{\arraystretch}{1.25}
\begin{tabularx}{\textwidth}{l X X}
\toprule
\textbf{Característica} & \textbf{Flannel (Práctica 3.1)} & \textbf{Calico (Práctica 1.4 y Trabajo Final)} \\
\midrule
Conexión local en nodo & Par \texttt{veth} conectado a un switch L2 (\texttt{cni0}, \texttt{10.244.X.1/24}). & Enlace punto a punto L3 \texttt{cali*} con ruta \texttt{/32} por Pod y \textit{proxy ARP} (\texttt{169.254.1.1}). \\
Transporte entre nodos & Túnel L2 sobre UDP: \textbf{VXLAN} (\texttt{flannel.1}, puerto UDP 8472, sobrecarga 50 bytes). & Enrutamiento L3 dinámico (\textbf{BGP} / BIRD) con encapsulación \textbf{IP-in-IP} (\texttt{tunl0}, protocolo IP 4, sobrecarga 20 bytes). \\
Soporte \texttt{NetworkPolicy} & \textbf{Nulo}: ignora cualquier objeto \texttt{NetworkPolicy} creado en la API. & \textbf{Nativo completo}: el agente \textbf{Felix} compila las políticas en cadenas \texttt{iptables}/\texttt{ipset} de alto rendimiento. \\
\bottomrule
\end{tabularx}
\caption{Comparativa técnica entre los plugins CNI Flannel y Calico en el laboratorio.}
\label{tab:flannel-vs-calico}
\end{table}

En Calico ocurren dos mecanismos de red esenciales que conviene distinguir con precisión:
\begin{enumerate}
  \item \textbf{Encaminamiento intra-nodo mediante \textit{Proxy ARP}}: dentro de cada Pod, la tabla de rutas presenta una entrada de enlace local (\texttt{169.254.1.1 dev eth0 scope link}) y la puerta de enlace por defecto vía \texttt{169.254.1.1}. Ninguna interfaz posee realmente esa IP: el kernel del nodo \textit{worker} activa \texttt{proxy\_arp=1} sobre la interfaz \texttt{cali*} correspondiente, respondiendo con su propia dirección MAC a cualquier solicitud ARP del Pod y encaminando inmediatamente el datagrama IP en Capa 3 según la tabla FIB del nodo, previa inspección en las cadenas \texttt{cali-fw-*} (\textit{from-workload}) y \texttt{cali-tw-*} (\textit{to-workload}) programadas por Felix.
  \item \textbf{Comunicación inter-nodo mediante túneles IP-in-IP (\texttt{tunl0})}: tal como se estudia en la Práctica 1.4 (Apdo.~3.3.1), cuando un Pod alojado en \texttt{k8s-worker1} envía un paquete por \texttt{eth0} a otro Pod alojado en \texttt{k8s-worker2} (p.\,ej., el Pod del controlador \texttt{ingress-nginx} alcanzando una réplica web remota), el demonio BGP (\texttt{BIRD}) de cada nodo anuncia su bloque CIDR local (\texttt{/26}) apuntando a la interfaz de túnel \textbf{\texttt{tunl0}}. El kernel encapsula el paquete IP original (\texttt{IP\_interna[SRC=PodA, DST=PodB]}) directamente como carga útil de un nuevo paquete IPv4 externo (\texttt{IP\_externa[SRC=10.10.10.11, DST=10.10.10.12, Protocol=4 (IPIP)]}) a través de \texttt{e1} (\texttt{k8s\_LAN}), ajustando la MTU de \texttt{eth0} en los Pods a \texttt{1480} bytes ($1500 - 20$ bytes de cabecera IPv4 exterior). Al capturar con Wireshark o \texttt{tcpdump} sobre \texttt{e1}, se observan limpiamente ambas cabeceras IP anidadas.
\end{enumerate}

\section{Meta-Plugin Multus y \texttt{NetworkAttachmentDefinition}}

La especificación CNI estándar invoca un único binario al crear el \texttt{netns} de un Pod. \textbf{Multus CNI} actúa como un \textit{meta-plugin} (un multiplexor CNI): es el primer binario que invoca \texttt{containerd} al leer \texttt{/etc/cni/net.d/00-multus.conf}. Su flujo de ejecución consta de dos pasos:
\begin{enumerate}
  \item Invoca siempre al \textbf{CNI primario} del clúster (\textbf{Calico}) para configurar \texttt{eth0}, la ruta por defecto del Pod y su registro en Kubernetes.
  \item Inspecciona las anotaciones (\path{metadata.annotations}) del Pod buscando la clave \path{k8s.v1.cni.cncf.io/networks}. Por cada entrada listada, consulta en el API Server el objeto CRD (\textit{Custom Resource Definition}) de tipo \path{NetworkAttachmentDefinition} correspondiente e invoca al plugin CNI secundario definido en su JSON (en nuestro caso, \texttt{macvlan} sobre \texttt{e2.1000}), inyectando la interfaz \texttt{net1} dentro del mismo \texttt{netns}.
\end{enumerate}

\section{Gestión de Cargas, Volúmenes, \texttt{ConfigMaps} y \texttt{Secrets}}

\subsection{\texttt{Deployment} frente a \texttt{Pod} Individual}

En nuestra arquitectura se emplean deliberadamente ambos niveles de abstracción según la naturaleza de cada componente:
\begin{itemize}
  \item \textbf{\texttt{Deployment} para los servidores web (\texttt{web-deployment})}: al tratarse de instancias sin estado propio (\textit{stateless}), un \texttt{Deployment} gestiona un \texttt{ReplicaSet} que mantiene de forma continua 3 réplicas activas (escalables en caliente cambiando \texttt{var.web\_replicas}). Si un Pod muere o un nodo cae, el controlador instancia otro Pod nuevo de forma autónoma (\textit{self-healing}).
  \item \textbf{\texttt{Pod} directo para el servidor \texttt{admin}}: al actuar como bastión único de gestión con IP fija (\texttt{10.1.2.10}) en la VLAN y anclado a \texttt{k8s-worker1} (\texttt{nodeSelector}) para recibir el tráfico SSH externo sin SNAT, se instancia como un recurso \texttt{Pod} individual controlado directamente por el estado de Terraform.
\end{itemize}

\subsection{Comparativa de Patrones: \texttt{initContainer} (Opción B) vs \texttt{sidecar} (Opción A)}

El Apartado 3 plantea personalizar el archivo \texttt{index.html} de cada réplica Nginx sustituyendo un marcador de plantilla por el nombre real del Pod. Evaluamos ambas arquitecturas:
\begin{itemize}
  \item \textbf{Opción A --- Contenedor auxiliar (\textit{sidecar})}: corre en paralelo junto a Nginx durante toda la vida del Pod. Resulta idóneo cuando el contenido cambia periódicamente (p.\,ej., sincronizando un repositorio Git cada 30 segundos), pero introduce dos inconvenientes para datos estáticos del Pod: mantiene un proceso extra consumiendo memoria/PIDs permanentemente, y genera una \emph{condición de carrera} durante el primer segundo de arranque en la que Nginx podría atender una petición antes de que el \textit{sidecar} haya terminado de escribir \texttt{index.html}.
  \item \textbf{Opción B --- \texttt{initContainer} (Solución elegida)}: los \texttt{initContainers} se ejecutan secuencialmente hasta finalizar con código \texttt{0} \emph{antes} de que el \texttt{kubelet} inicie los contenedores principales de aplicación. Dado que \texttt{metadata.name}, \texttt{spec.nodeName} y \texttt{status.podIP} son constantes durante la vida del Pod (expuestos de forma nativa mediante variables de entorno de la \textbf{\textit{Downward API}}), el \texttt{initContainer} \texttt{generate-html} lee la plantilla desde el volumen \texttt{ConfigMap} (montado como solo lectura en \texttt{/mnt/html-template}), ejecuta \texttt{sed}, deposita el \texttt{index.html} resultante en el volumen \texttt{emptyDir} compartido (\texttt{/usr/share/nginx/html}) y termina limpiamente liberando el 100\,\% de sus recursos antes de que arranque Nginx.
\end{itemize}

\subsection{Montaje Seguro de \texttt{Secrets} y Propiedad \texttt{fsGroup}}

Un \texttt{Secret} de Kubernetes montado como volumen utiliza un sistema de ficheros temporal en RAM (\texttt{tmpfs}) para evitar que las claves toquen el disco físico del nodo \textit{worker}. Por defecto, \texttt{kubelet} crea los ficheros del volumen con propietario \texttt{root:root} (\texttt{0:0}) y permisos \texttt{0644}. Al especificar en el Pod:
\begin{lstlisting}[language=YAML]
spec:
  securityContext:
    fsGroup: 1000
  volumes:
    - name: ssh-private-key
      secret:
        secretName: admin-ssh-privkey
        defaultMode: 0440
\end{lstlisting}
Kubernetes asigna la propiedad de grupo \texttt{GID=1000} (\texttt{adminuser}/\texttt{webuser}) a los puntos de montaje y aplica la máscara octal \texttt{0440} (\texttt{-r--r-----}, propietario \texttt{root}, grupo \texttt{1000}), permitiendo que el usuario sin privilegios de cada contenedor lea el secreto de manera segura a través de sus permisos de grupo, sin exponerlo a otros usuarios (\textit{others = 0}).

\section{Servicios, MetalLB, Ingress NGINX y Modelo Stateful de \texttt{NetworkPolicy}}

La publicación y protección de los servicios combina tres capas sucesivas:
\begin{enumerate}
  \item \textbf{MetalLB (Capa 2 / ARP)}: en ausencia de un balanceador de proveedor cloud (AWS ELB, GCP LB), el demonio \texttt{speaker} de MetalLB anuncia mediante respuestas ARP gratuitas en \texttt{k8s\_LAN} (\texttt{e1}) la IP externa \texttt{10.10.10.20} asignada al servicio \texttt{LoadBalancer} del controlador \texttt{ingress-nginx}.
  \item \textbf{Ingress NGINX (Capa 7 HTTP)}: un único punto de entrada (\texttt{10.10.10.20:80}) inspecciona la URL HTTP entrante. Cuando coincide con el prefijo \texttt{/servers-G11}, la anotación \texttt{nginx.ingress.kubernetes.io/rewrite-target: "/"} reescribe la ruta URI a \texttt{"/"} (pues Nginx en los Pods sirve su raíz \texttt{/usr/share/nginx/html/index.html}) y encamina la petición hacia los \textit{Endpoints} de \texttt{web-service} (\texttt{ClusterIP:8080 $\rightarrow$ PodIP:80}).
  \item \textbf{Modelo aditivo y \textit{stateful} de \texttt{NetworkPolicy} (Práctica 1.4, Apdo.~3.3)}: en Kubernetes, si ninguna política selecciona a un Pod, todo su tráfico está permitido (\textit{default allow}). Al aplicar nuestra primera política \texttt{default-deny-all} (\texttt{podSelector: \{\}} con \texttt{policyTypes: ["Ingress", "Egress"]} y sin reglas \texttt{ingress}/\texttt{egress}), todo el tráfico entrante y saliente en el namespace \texttt{default} queda denegado salvo lo permitido explícitamente por otras políticas (modelo aditivo o \textit{whitelist}).

  Tal como demuestra el estudio experimental de la Práctica 1.4 (Apdo.~3.3.3), Calico evalúa las políticas de manera \textbf{con estado (\textit{stateful})} apoyándose en \texttt{conntrack}: distingue quién \emph{inicia} la conexión (\texttt{TCP SYN}) del tráfico de \emph{respuesta}. Cuando una comunicación ocurre entre dos Pods sujetos ambos a \texttt{default-deny-all} dentro del mismo namespace (como el Pod \texttt{tester} y los Pods \texttt{nginx} en la Práctica 1.4), es obligatorio definir \textbf{tanto} una regla \texttt{Egress} en el Pod emisor \textbf{como} una regla \texttt{Ingress} en el Pod receptor para autorizar el \texttt{SYN} inicial, mientras que los paquetes de respuesta regresan automáticamente sin reglas inversas. En cambio, en nuestro Trabajo Final, tanto el controlador \texttt{ingress-nginx} (que reside en su propio namespace \texttt{ingress-nginx}, no afectado por \texttt{default-deny-all}) como el host SSH exterior (\texttt{10.10.20.2}) se encuentran \emph{fuera} del namespace \texttt{default}: por ello, basta definir únicamente reglas \texttt{Ingress} (\texttt{allow-ingress-to-web} y \texttt{allow-ssh-to-admin}) en nuestros Pods sin abrir ninguna regla \texttt{Egress} en \texttt{eth0}, ya que todo el tráfico HTTP y SSH de vuelta fluye autorizado por el seguimiento de estado de la conexión entrante.
\end{enumerate}

\begin{warnbox}{Semántica Booleana AND vs OR en \texttt{NetworkPolicy}}
En la política \texttt{allow-ingress-to-web}, colocar \texttt{namespaceSelector} y \texttt{podSelector} dentro de \textbf{un mismo bloque \texttt{from}} impone una condición \textbf{AND lógica} (únicamente los Pods con etiquetas de \texttt{ingress-nginx} que residan dentro del namespace \texttt{ingress-nginx}). Si se declarasen como dos bloques \texttt{from} separados, Kubernetes aplicaría un \textbf{OR lógico}, abriendo una brecha de seguridad que permitiría el acceso desde cualquier Pod del namespace \texttt{ingress-nginx} o desde cualquier Pod del namespace \texttt{default} que falsificase esas etiquetas.
\end{warnbox}

\chapter{Orquestación Declarativa con Terraform: Fundamentos e Implementación}

Este capítulo sintetiza de forma completa el funcionamiento de \textbf{Terraform} ---cubriendo tanto todos los conceptos, bloques HCL y patrones introducidos en la \textbf{Práctica 1.4} como las funcionalidades avanzadas aplicadas en el \textbf{Trabajo Final}--- para servir como referencia integral de estudio y documentación técnica del despliegue.

\section{Filosofía IaC y Arquitectura del Nodo \texttt{terraform} en el Laboratorio}

En las prácticas iniciales (\texttt{p21} sobre OpenStack CLI y \texttt{p31} sobre \texttt{kubectl}), la creación del escenario se realizaba mediante secuencias imperativas de comandos (\texttt{openstack server create}, \texttt{kubectl create}, \texttt{kubectl apply}). El enfoque imperativo presenta limitaciones importantes: si un script se interrumpe a mitad de ejecución, relanzarlo falla por colisión de nombres; modificar parámetros exige comandos manuales de parcheo; y no existe un registro único de recursos al desmontar la infraestructura.

\textbf{Terraform} (HashiCorp) aplica el paradigma declarativo de \textbf{Infraestructura como Código (IaC)} mediante el lenguaje \textbf{HCL (\textit{HashiCorp Configuration Language})}, ofreciendo un motor uniforme multi-proveedor (capaz de orquestar tanto máquinas virtuales, redes, routers e IPs flotantes en \textbf{OpenStack} con \path{terraform-provider-openstack/openstack} como cargas de contenedores y políticas en \textbf{Kubernetes} con \path{hashicorp/kubernetes}):
\begin{itemize}
  \item \textbf{Idempotencia}: aplicar reiteradamente una misma plantilla HCL cuando el clúster ya converge con ella produce cero operaciones (\texttt{0 added, 0 changed, 0 destroyed}).
  \item \textbf{Reconciliación a tres bandas y detección de deriva (\textit{drift})}: Terraform contrasta la especificación deseada (\texttt{*.tf}), el estado registrado (\texttt{terraform.tfstate}) y el estado real consultado en vivo a la API del proveedor. Si un recurso se altera manualmente fuera de Terraform o se incrementa \texttt{web\_replicas} de 3 a 5, \texttt{terraform plan} calcula el diferencial exacto sin destruir el resto de la arquitectura.
  \item \textbf{Gestión automática del ciclo de vida}: ante cambios compatibles (p.\,ej., escalar réplicas de un \texttt{Deployment}), Terraform actualiza el objeto \textit{in-place}; ante cambios en campos inmutables (p.\,ej., cambiar el \texttt{nodeSelector} o nombre de un \texttt{Pod} directo o el \textit{keypair} de una VM en OpenStack), destruye y recrea automáticamente el recurso afectado.
\end{itemize}

\subsection{Integración del Nodo Proveedor \texttt{terraform} en VNX (Práctica 1.4)}

En el entorno docente de DIT-UPM (\path{/mnt/tmp/k8s_lab-calico-v02}), el cliente de Terraform no reside en el sistema anfitrión, sino en un nodo VNX dedicado llamado \textbf{\texttt{terraform}} (\texttt{10.10.10.13} en \texttt{k8s\_LAN}, con salida NAT por \texttt{virbr0} para descargar plugins desde \path{registry.terraform.io}). Nuestro script \texttt{01\_deploy\_all.sh} automatiza la secuencia exacta de puesta en marcha definida en la Práctica 1.4:
\begin{lstlisting}[language=bash,caption={Aprovisionamiento del nodo \texttt{terraform} y exportación de credenciales del clúster.}]
# 1. Arranque del cluster Kubernetes con Calico
/lab/cnvr/bin/get-k8s-calico
cd /mnt/tmp/k8s_lab-calico-v02
./bin/start-k8s-calico

# 2. Creacion del nodo VNX 'terraform' (10.10.10.13) e instalacion del binario CLI
sudo vnx -f k8s_lab-terraform.xml --create
sudo vnx -f k8s_lab-terraform.xml -x install-terraform

# 3. Carga de kubeconfig en el host y exportacion a /root/.kube/config en 'terraform'
./bin/get-master-kubeconf
./bin/set-terraform-kubeconf
\end{lstlisting}
Entre el host (\path{/mnt/tmp/k8s_lab-calico-v02/shared}) y la máquina \texttt{terraform} (\texttt{/shared}) existe un sistema de ficheros compartido que permite editar o sincronizar los archivos \texttt{*.tf} desde el host y ejecutarlos dentro de \texttt{terraform}. Dado que \texttt{/mnt/tmp} es volátil, nuestro proyecto conserva todo el código fuente en su directorio persistente y respalda automáticamente el \texttt{terraform.tfstate} generado tras cada despliegue.

\section{Bloques Fundamentales de HCL y Ciclo de Comandos}

\subsection{Flujo Operativo de Comandos y Operación Selectiva con \texttt{-target}}

La operación con Terraform sigue un ciclo estructurado de seis comandos principales:
\begin{enumerate}
  \item \textbf{\texttt{terraform init}} (y \textbf{\texttt{terraform init -upgrade}} al añadir nuevos módulos o actualizar versiones): inicializa el directorio de trabajo, descarga los binarios de los proveedores declarados en \texttt{.terraform/} y registra sus firmas criptográficas en \path{.terraform.lock.hcl}.
  \item \textbf{\texttt{terraform validate}}: comprueba localmente la corrección sintáctica HCL, la coherencia de tipos de variables y la validez de atributos sin contactar aún con la API remota.
  \item \textbf{\texttt{terraform plan}}: consulta el estado real en el proveedor, lo compara con el código HCL y muestra la vista previa de recursos a crear (\texttt{+}), modificar \textit{in-place} (\texttt{\textasciitilde}) o destruir/recrear (\texttt{-/+}), pudiendo guardar el plan firmado con \texttt{-out=tfplan}.
  \item \textbf{\texttt{terraform apply}}: ejecuta los cambios planificados y actualiza \texttt{terraform.tfstate}.
  \item \textbf{\texttt{terraform state list}} y \textbf{\texttt{terraform state show <recurso>}}: permiten auditar el inventario completo de objetos bajo control de Terraform y examinar todos los atributos en vivo de un recurso concreto (p.\,ej., \texttt{terraform state show} \path{kubernetes_deployment_v1.web}).
  \item \textbf{\texttt{terraform destroy}} y el modificador \textbf{\texttt{-target}}: mientras que \texttt{terraform destroy} elimina ordenadamente toda la infraestructura, añadir \path{-target=<tipo>.<nombre>} (tanto en \texttt{plan}, \texttt{apply} como \texttt{destroy}) restringe la acción exclusivamente a un subgrafo concreto (p.\,ej., \texttt{terraform apply} \path{-target=kubernetes_network_policy_v1.default_deny_all}), muy útil durante pruebas incrementales de políticas de seguridad.
\end{enumerate}

\subsection{Tipos de Bloques HCL: \texttt{provider}, \texttt{variable}, \texttt{tfvars}, \texttt{data}, \texttt{output} y \texttt{module}}

\begin{itemize}
  \item \textbf{Bloques \texttt{terraform} y \texttt{provider} (\path{k8s-provider.tf})}: siguiendo la convención exacta de las plantillas de la Práctica 1.4 (\path{apartado3-1/k8s-provider.tf}), fija \texttt{required\_version = ">= 0.14.0"} y la versión exacta \texttt{version = "2.23.0"} de \path{hashicorp/kubernetes} (garantizando compatibilidad total con el binario instalado por \texttt{install-terraform} en la VM del laboratorio) e indica \texttt{config\_path = pathexpand(var.kubeconfig\_path)}.
  \item \textbf{Variables de entrada (\texttt{variables.tf}) y los 3 métodos de parametrización}: cada bloque \texttt{variable} define \texttt{description}, \texttt{type} (\texttt{string}, \texttt{number}, \texttt{list(string)}, \texttt{map(string)}), valor \texttt{default} y reglas \texttt{validation}. Sus valores pueden suministrarse por tres vías (en orden creciente de precedencia sobre \texttt{default}):
    \begin{enumerate}
      \item \textbf{Carga automática mediante \texttt{terraform.tfvars}}: si el directorio contiene un fichero llamado exactamente \path{terraform.tfvars} (incluido en nuestro repositorio particularizado para el Grupo 11), Terraform lo carga automáticamente en cada \texttt{plan}/\texttt{apply} sin banderas adicionales.
      \item \textbf{Fichero explícito}: \texttt{terraform apply} \path{-var-file=parametros.tfvars}.
      \item \textbf{Línea de comandos (máxima precedencia)}: \texttt{terraform apply} \path{-var="web_replicas=5"}, empleada por \path{01_deploy_all.sh} para inyectar parámetros autodetectados en tiempo de ejecución.
    \end{enumerate}
  \item \textbf{Fuentes de datos (\texttt{data}) y salidas (\texttt{outputs.tf})}: a diferencia de \texttt{resource} (que crea y gobierna un objeto), un bloque \textbf{\texttt{data}} consulta objetos en solo lectura. En \texttt{terraform/outputs.tf} aplicamos el patrón de consulta plural \texttt{data "kubernetes\_resources"} combinado con una expresión de comprensión de listas \texttt{for} (Práctica 1.4, Apdo.~3.1) para devolver dinámicamente los nombres exactos de los Pods generados por el \texttt{Deployment} tras su creación:
\begin{lstlisting}[language=HCL,caption={Consulta dinámica de Pods con \texttt{data "kubernetes\_resources"} y bloque \texttt{output} en \texttt{outputs.tf}.}]
data "kubernetes_resources" "web_pods" {
  api_version    = "v1"
  kind           = "Pod"
  namespace      = var.namespace
  label_selector = "app=${local.web_labels.app}"
  depends_on     = [kubernetes_deployment_v1.web]
}

output "web_pod_names" {
  value = [for pod in data.kubernetes_resources.web_pods.objects : pod.metadata.name]
}
\end{lstlisting}
  Los valores expuestos pueden consultarse en cualquier momento con \texttt{terraform output}, de forma individual (\texttt{terraform output web\_ingress\_url}) o en formato estructurado para scripts (\texttt{terraform output -json}).
  \item \textbf{Módulos reutilizables (\texttt{module})}: cuando un patrón de recursos se instancia múltiples veces con distintos parámetros (como el módulo \texttt{modules/vm\_instance/\{main.tf, variables.tf, outputs.tf\}} de la Práctica 1.4 Apdo.~2.3), Terraform permite encapsularlo e invocarlo mediante bloques \texttt{module "<nombre>" \{ source = "./modules/..." \}}, accediendo a sus resultados como \texttt{module.<nombre>.<salida>}.
\end{itemize}

\section{Las Cuatro Formas de Referenciar Recursos de Kubernetes en Terraform}

Un aspecto clave analizado en la Práctica 1.4 (Apdo.~3.2) es cómo vincular el \texttt{selector} de un \texttt{Service} o \texttt{NetworkPolicy} con las etiquetas (\texttt{labels}) de los Pods de un \texttt{Deployment}. La Tabla~\ref{tab:cuatro-referencias} compara las cuatro alternativas existentes en HCL:

\begin{table}[H]
\centering
\footnotesize
\renewcommand{\arraystretch}{1.25}
\begin{tabularx}{\textwidth}{l >{\raggedright\arraybackslash}p{4.8cm} >{\raggedright\arraybackslash}X}
\toprule
\textbf{Estrategia} & \textbf{Expresión HCL en \texttt{selector}} & \textbf{Ventajas, Limitaciones y Uso} \\
\midrule
1. Etiquetas literales & \texttt{selector = \{ app = "web" \}} & Sencilla y directa (idéntica a YAML), pero propensa a desincronización humana si cambian las etiquetas y \textbf{no informa a Terraform} de dependencia en el DAG. \\
2. Referencia directa a recurso & \path{kubernetes_deployment_v1.web.spec[0].template[0].metadata[0].labels} & Crea \textbf{dependencia implícita automática} en el grafo DAG y propaga cambios de etiquetas; requiere indexar con \texttt{[0]} cada sub-bloque del esquema y resulta verbosa. \\
3. Fuente de datos \texttt{data} & \path{data.kubernetes_resource.web.object.spec.template.metadata.labels} & Opción idónea cuando el \texttt{Deployment} objetivo \textbf{ya existía previamente} en el clúster fuera del control de nuestra plantilla Terraform. \\
4. Valores locales (\texttt{locals}) & \texttt{selector = local.web\_labels} (definido en \texttt{locals.tf}) & \textbf{Fuente única de verdad} limpia y legible compartida entre \texttt{Deployment}, \texttt{Service} y \texttt{NetworkPolicy}; al no apuntar de un recurso a otro, no añade arista implícita por sí sola. \\
\bottomrule
\end{tabularx}
\caption{Comparativa de los cuatro mecanismos de referencia entre recursos Kubernetes en Terraform.}
\label{tab:cuatro-referencias}
\end{table}

En nuestra arquitectura (\texttt{terraform/}) combinamos lo mejor de los métodos \textbf{4} y \textbf{2}: centralizamos las etiquetas en \textbf{\path{locals.tf}} (\path{local.web_labels} y \path{local.admin_labels}) para mantener una legibilidad impecable sin duplicación en \path{web.tf}, \path{admin.tf}, \path{services_ingress.tf}, \path{network_policies.tf} y \path{outputs.tf}, mientras que enlazamos nombres de objetos mediante referencias directas (\path{kubernetes_service_v1.web.metadata[0].name}, \path{kubernetes_secret_v1.*.metadata[0].name}) y sentencias explícitas \texttt{depends\_on} donde la dependencia opera en el plano de datos CNI.

\section{Grafo de Dependencias (DAG), \texttt{kubernetes\_manifest} y Sutilezas HCL2}

Durante \texttt{plan} y \texttt{apply}, Terraform construye un \textbf{Grafo Acíclico Dirigido (DAG)} para paralelizar la creación de objetos independientes respetando el orden estricto entre capas (Figura~\ref{fig:dag-terraform}):

\begin{figure}[H]
\centering
\begin{tikzpicture}[
  font=\sffamily\scriptsize,
  tfnode/.style={draw=upmblue, rounded corners=3pt, fill=white, thick, align=center, inner sep=4.5pt, minimum width=3.2cm},
  crdnode/.style={draw=upmaccent, rounded corners=3pt, fill=upmaccent!8, thick, align=center, inner sep=4.5pt, minimum width=3.2cm},
  npnode/.style={draw=errred!80!black, rounded corners=3pt, fill=errred!6, thick, align=center, inner sep=4.5pt, minimum width=3.2cm},
  datanode/.style={draw=okgreen!80!black, rounded corners=3pt, fill=okgreen!7, thick, align=center, inner sep=4.5pt, minimum width=3.2cm},
  dep/.style={-{Stealth[length=2mm]}, thick, color=darkslate},
  expdep/.style={-{Stealth[length=2mm]}, thick, dashed, color=upmaccent}
]

  % Nivel 1: Recursos base paralelos
  \node[tfnode] (sec) at (-4.5, 2.6) {\textbf{\texttt{secrets.tf}}\\3 $\times$ \texttt{kubernetes\_secret\_v1}};
  \node[tfnode] (cm)  at (0.0, 2.6)  {\textbf{\texttt{configmap.tf}}\\\texttt{web-html-template}};
  \node[crdnode] (nad) at (4.5, 2.6) {\textbf{\texttt{multus.tf}}\\2 $\times$ \texttt{kubernetes\_manifest} (NAD)};

  % Nivel 2: Workloads y Services
  \node[tfnode] (adm) at (-4.5, 0.7) {\textbf{\texttt{admin.tf}}\\\texttt{kubernetes\_pod\_v1.admin}};
  \node[tfnode] (web) at (0.0, 0.7)  {\textbf{\texttt{web.tf}}\\\texttt{kubernetes\_deployment\_v1.web}};
  \node[tfnode] (ic)  at (4.5, 0.7)  {\textbf{\texttt{services\_ingress.tf}}\\\texttt{IngressClass} + \texttt{Services}};

  % Nivel 3: Policies, Data Pod List e Ingress
  \node[npnode]   (np)  at (-4.5, -1.3) {\textbf{\texttt{network\_policies.tf}}\\3 $\times$ \texttt{NetworkPolicy}};
  \node[datanode] (dat) at (0.0, -1.3)  {\textbf{\texttt{outputs.tf}}\\\texttt{data.kubernetes\_resources}};
  \node[tfnode]   (ing) at (4.5, -1.3)  {\textbf{\texttt{services\_ingress.tf}}\\\texttt{kubernetes\_ingress\_v1.web}};

  % Aristas implicitas (solidas) y explicitas depends_on (discontinuas)
  \draw[dep] (sec) -- (adm);
  \draw[dep] (sec.south east) -- (web.north west);
  \draw[dep] (cm)  -- node[right, font=\sffamily\tiny] {\texttt{sha256()}} (web);
  \draw[expdep] (nad.south west) -- node[pos=0.5, above, sloped, font=\sffamily\tiny] {\texttt{depends\_on}} (web.north east);
  \draw[expdep] (nad.south) to[out=-145,in=20] (adm.north east);

  \draw[expdep] (adm) -- (np);
  \draw[expdep] (web.south west) -- (np.north east);
  \draw[expdep] (web) -- (dat);
  \draw[dep] (ic) -- (ing);

\end{tikzpicture}
\caption{Grafo de dependencias (DAG) en Terraform: líneas continuas para referencias implícitas HCL y discontinuas para dependencias explícitas \texttt{depends\_on}.}
\label{fig:dag-terraform}
\end{figure}

Obsérvese por qué \texttt{depends\_on} en \texttt{web.tf} y \texttt{admin.tf} hacia \texttt{kubernetes\_manifest.vlan\_*} es crítico: para Kubernetes, el nombre de la red Multus es una cadena de texto en una anotación (\texttt{local.vlan\_dhcp\_nad\_name}). Sin \texttt{depends\_on}, Terraform lanzaría en paralelo la creación de la \texttt{NetworkAttachmentDefinition} y de los Pods, provocando errores \texttt{FailedCreatePodSandBox} si el \texttt{kubelet} consultase el CRD milisegundos antes de su registro en \texttt{etcd}.

\subsection{CRDs Personalizados con \texttt{kubernetes\_manifest} y Sutilezas HCL2}

Mientras que los objetos estándar usan recursos tipados nativos (\texttt{kubernetes\_deployment\_v1}, \texttt{kubernetes\_secret\_v1}), los Custom Resource Definitions de terceros como \path{NetworkAttachmentDefinition} (\path{k8s.cni.cncf.io/v1}) requieren el recurso genérico \textbf{\texttt{kubernetes\_manifest}} (Práctica 1.4, Apdo.~3.4), que emplea \textit{Server-Side Apply} validando el objeto contra el esquema OpenAPI del \texttt{kube-apiserver}:

\begin{lstlisting}[language=HCL,caption={Definición declarativa de las redes Multus en \texttt{terraform/multus.tf} usando \texttt{locals} y \texttt{jsonencode()}.}]
resource "kubernetes_manifest" "vlan_dhcp" {
  manifest = {
    apiVersion = "k8s.cni.cncf.io/v1"
    kind       = "NetworkAttachmentDefinition"
    metadata = {
      name      = local.vlan_dhcp_nad_name
      namespace = var.namespace
    }
    spec = {
      config = jsonencode({
        cniVersion = "0.3.1"
        name       = local.vlan_dhcp_nad_name
        type       = "macvlan"
        master     = "${var.worker_vlan_parent_iface}.${var.vlan_id}"
        mode       = "bridge"
        ipam       = { type = "dhcp" }
      })
    }
  }
  field_manager { force_conflicts = true }
}
\end{lstlisting}

Finalmente, destacan cuatro detalles técnicos de HCL2 aplicados en nuestros módulos:
\begin{enumerate}
  \item \textbf{Inyección nativa de claves SSH con \texttt{file()}}: en \texttt{secrets.tf}, \texttt{kubernetes\_secret\_v1} lee directamente las claves generadas en \texttt{../keys/} mediante \texttt{file("\$\{path.module\}/...")}, codificándolas automáticamente en Base64 hacia la API de Kubernetes sin usar \texttt{kubectl}.
  \item \textbf{Permisos octales como cadenas (\texttt{"0440"})}: HCL2 no interpreta literales numéricos con cero inicial como octales (\texttt{0440} numérico sería decimal \texttt{440}). Pasar \texttt{default\_mode = "0440"} y \texttt{mode = "0440"} como cadenas en \texttt{admin.tf} y \texttt{web.tf} aplica el modo octal exacto (\texttt{-r--r-----}).
  \item \textbf{Escape de variables shell \texttt{\$\$} en bloques \textit{heredoc}}: dentro del comando del \texttt{initContainer} (\texttt{<<-EOT}), escribir \texttt{\$\$\{POD\_NAME\}} escapa la interpolación de Terraform para que llegue literalmente como \texttt{\$\{POD\_NAME\}} al intérprete \texttt{/bin/sh} del contenedor.
  \item \textbf{Actualización rodante automática con \texttt{sha256()}}: la anotación \path{"cnvr.upm.es/template-hash"} calculada con \path{sha256(kubernetes_config_map_v1.web_html_template.data["index.html"])} provoca un \textit{rolling update} automático de las réplicas web cada vez que se edita el HTML del \texttt{ConfigMap}.
\end{enumerate}

\chapter{Decisiones de Ingeniería y Resultados (\textit{Decisions \& Outcomes})}

Este capítulo documenta de forma razonada los seis problemas reales de integración entre Linux, VNX, Kubernetes y Terraform encontrados al diseñar el escenario, las opciones evaluadas y el resultado técnico (\textit{outcome}) de cada decisión adoptada.

\section{Decisión 1: Asignación de IPs en Multus para un \texttt{Deployment} Escalable}

\begin{decisionbox}{Problema: Colisión de IPs en Capa 2 al escalar réplicas con Multus}
En la Práctica 3.1 (\texttt{13\_multus.sh}), cada Pod se definía en un fichero individual (\texttt{vma3.yml}, \texttt{vma4.yml}) con una anotación \texttt{ips: ["10.1.2.200/24"]} bajo \texttt{ipam: static}. Sin embargo, en el Trabajo Final los tres servidores web pertenecen a un único objeto \textbf{\texttt{Deployment}}, donde todas las réplicas comparten exactamente la misma plantilla (\texttt{spec.template}).
\begin{itemize}
  \item Si se fija una IP estática en la plantilla, las 3 réplicas arrancan con la \emph{misma IP} en la VLAN \texttt{1000}, causando conflicto ARP inmediato.
  \item Si se utiliza el plugin \texttt{host-local} con un mismo rango en la \texttt{NetworkAttachmentDefinition}, como su base de datos reside localmente en el disco de cada nodo (\texttt{/var/lib/cni/networks}), \texttt{k8s-worker1} asignaría \texttt{10.1.2.50} a su primer Pod y \texttt{k8s-worker2} asignaría \emph{también} \texttt{10.1.2.50} al suyo.
\end{itemize}
\end{decisionbox}

\textbf{Solución de ingeniería implementada:}
Se desacopla la política IPAM en dos definiciones \texttt{NetworkAttachmentDefinition} complementarias (\texttt{multus.tf}):
\begin{enumerate}
  \item \textbf{\texttt{vlan1000-static}} para el Pod único \texttt{admin}, fijando su dirección conocida en \texttt{10.1.2.10/24}.
  \item \textbf{\texttt{vlan1000-dhcp}} para el \texttt{Deployment} web: en \texttt{01\_deploy\_all.sh} se activa el demonio oficial \texttt{/opt/cni/bin/dhcp daemon} (\texttt{/run/cni/dhcp.sock}) en ambos nodos \textit{worker}, mientras que el contenedor LXC \texttt{logs} (\texttt{10.1.2.101}) ejecuta \textbf{\texttt{dnsmasq}} sobre \texttt{eth1} repartiendo concesiones únicas del rango \texttt{10.1.2.50--10.1.2.95}. Crucialmente, en \texttt{vnx/conf/dnsmasq/vlan1000-dhcp.conf} se configura \texttt{port=0}, \texttt{dhcp-option=3} (vacío) y \texttt{dhcp-option=6} (vacío) para que el servidor DHCP entregue \textbf{únicamente la dirección IP \texttt{/24} de enlace}, sin inyectar una segunda puerta de enlace por defecto que sobreescriba la ruta primaria de Calico en \texttt{eth0} ni alterar \texttt{/etc/resolv.conf} (CoreDNS).
\end{enumerate}

\begin{outcomebox}{Resultado (\textit{Outcome})}
Cualquier operación \texttt{terraform apply -var="web\_replicas=N"} (con $N = 3, 5, 10\dots$) asigna direcciones IP únicas y libres de colisiones en toda la VLAN \texttt{1000} con independencia del nodo \textit{worker} elegido por el \texttt{kube-scheduler}, manteniendo intacto el encaminamiento primario de Kubernetes.
\end{outcomebox}

\section{Decisión 2: \texttt{Secrets}, \texttt{fsGroup: 1000} y \texttt{StrictModes} de OpenSSH}

\begin{decisionbox}{Problema: Rechazo de claves privadas \texttt{0440} por el cliente OpenSSH}
El enunciado propone montar los recursos \texttt{Secret} con \texttt{defaultMode: 0440} y \texttt{securityContext.fsGroup: 1000} para que \texttt{adminuser} (\texttt{UID=1000, GID=1000}) pueda leer el archivo montado (propiedad \texttt{root:1000}, permisos \texttt{-r--r-----}). No obstante, al inspeccionar el código fuente de OpenSSH (\texttt{sshkey\_perm\_ok}), el cliente \texttt{ssh} comprueba estrictamente que \texttt{(st\_uid == uid) \&\& ((st\_mode \& 077) == 0)}. Si un usuario no-root invoca \texttt{ssh -i} contra un fichero propiedad de \texttt{root} o con bits de grupo activos (\texttt{0440 \& 077 = 040 $\neq$ 0}), OpenSSH aborta la conexión con \texttt{WARNING: UNPROTECTED PRIVATE KEY FILE! Permissions 0440 are too open}.
\end{decisionbox}

\textbf{Solución de ingeniería implementada:}
\begin{itemize}
  \item Se generan en el host dos pares RSA de 4096 bits completamente independientes: \textbf{Par 1} (\texttt{keys/id\_rsa\_internal}, administración interna) y \textbf{Par 2} (\texttt{keys/id\_rsa\_external}, acceso exterior al bastión).
  \item En \texttt{terraform/secrets.tf} se declaran \textbf{tres objetos \texttt{kubernetes\_secret\_v1} separados}: \texttt{admin-ssh-privkey} (clave privada Par 1), \texttt{admin-ssh-authkeys} (clave pública Par 2) y \texttt{web-ssh-authkeys} (clave pública Par 1), montados todos ellos con \texttt{mode = "0440"} y \texttt{fs\_group = 1000}.
  \item En \path{admin.tf} y \path{web.tf} se añade un \texttt{initContainer} ejecutado explícitamente sin privilegios (\path{run_as_user = 1000}, \path{run_as_group = 1000}). Este contenedor \emph{depende} de \path{fsGroup: 1000} y del modo \texttt{0440} para poder leer los \texttt{Secrets} montados en \path{/mnt/ssh-*}, y copia las claves a un volumen \texttt{emptyDir} montado en \path{/home/adminuser/.ssh} (y \path{/home/webuser/.ssh}) aplicando directorio \texttt{0700} y ficheros \texttt{0600} propiedad de \texttt{1000:1000}. Adicionalmente, los volúmenes \texttt{Secret} originales permanecen montados en \path{/etc/ssh-secrets/} para su inspección directa.
\end{itemize}

\begin{outcomebox}{Resultado (\textit{Outcome})}
Se cumple al 100\,\% el diseño de \texttt{Secrets} + \texttt{fsGroup: 1000} + \texttt{0440} en Kubernetes y, simultáneamente, tanto el demonio \texttt{sshd} (\texttt{StrictModes yes}) como el cliente \texttt{ssh} autentican de forma inmediata y transparente sin contraseñas ni errores de permisos.
\end{outcomebox}

\section{Decisión 3: Inspección de Imágenes OCI y Renderizado HTML Dinámico}

Al inspeccionar los manifiestos OCI de \path{dgsgiros/nginx-ssh:latest} y \path{dgsgiros/admin-ssh:latest}, se verificó que:
\begin{itemize}
  \item \path{dgsgiros/nginx-ssh:latest} extiende la imagen oficial \texttt{nginx:1.31.8} sobre Debian Trixie, sirviendo contenido desde \textbf{\path{/usr/share/nginx/html}}, con un \path{/entrypoint.sh} que ejecuta \texttt{ssh-keygen -A}, lanza \path{/usr/sbin/sshd} en segundo plano y transfiere el PID 1 a \texttt{nginx -g 'daemon off;'}.
  \item Reutilizar esa misma imagen \texttt{dgsgiros/nginx-ssh:latest} en los \texttt{initContainers} de \texttt{web.tf} evita descargar imágenes extra (\texttt{busybox}/\texttt{alpine}) desde Docker Hub, acelerando el arranque de cada réplica y reduciendo dependencias de red.
\end{itemize}

\begin{lstlisting}[language=HCL,caption={Fragmento del \texttt{initContainer} (Opción B) en \texttt{terraform/web.tf} con \textit{Downward API}.}]
init_container {
  name  = "generate-html"
  image = "dgsgiros/nginx-ssh:latest"
  env {
    name = "POD_NAME"
    value_from { field_ref { field_path = "metadata.name" } }
  }
  command = ["/bin/sh", "-c", <<-EOT
    sed -e "s/__POD_NAME__/$${POD_NAME}/g" \
        -e "s/__GROUP_NUMBER__/$${GROUP_NUMBER}/g" \
        /mnt/html-template/index.html > /usr/share/nginx/html/index.html
  EOT
  ]
  volume_mount { name = "html-template-vol", mount_path = "/mnt/html-template", read_only = true }
  volume_mount { name = "html-shared-vol",   mount_path = "/usr/share/nginx/html" }
}
\end{lstlisting}

\begin{outcomebox}{Resultado (\textit{Outcome})}
Cada Pod sirve desde el primer milisegundo en que su \texttt{readinessProbe} pasa a \texttt{True} una tarjeta HTML personalizada con su nombre exacto de Pod (\texttt{web-deployment-<rs-hash>-<pod-id>}), el nodo \textit{worker} que lo ejecuta, su IP primaria en Calico y la ruta \texttt{/servers-G11}.
\end{outcomebox}

\section{Decisión 4: Preservación de IP Origen en \texttt{NodePort} para \texttt{ipBlock}}

\begin{decisionbox}{Problema: El SNAT por defecto de \texttt{kube-proxy} anula la \texttt{NetworkPolicy} de SSH}
El Apartado 5.3 exige permitir SSH al Pod \texttt{admin} exclusivamente desde la IP del host autorizado en \texttt{user\_LAN} (\texttt{10.10.20.2/32}) usando \texttt{ipBlock}. Si \texttt{admin-service} mantuviese \texttt{externalTrafficPolicy: Cluster}, el paquete procedente de \texttt{10.10.20.2} hacia \texttt{10.10.10.11:30022} recibiría SNAT en \texttt{k8s-worker1} pasando a tener \texttt{SRC=10.10.10.11}, por lo que Calico lo descartaría al no coincidir con \texttt{10.10.20.2/32}.
\end{decisionbox}

\textbf{Solución y recorrido verificado del paquete:}
En \texttt{admin.tf} se fija \path{node_selector = { "kubernetes.io/hostname" = "k8s-worker1" }} y en \path{services_ingress.tf} se configura \path{external_traffic_policy = "Local"} sobre \texttt{admin-service} (\texttt{port = 2222, target\_port = 22, node\_port = 30022}). El recorrido exacto del paquete es:
\begin{enumerate}
  \item El host emite \texttt{IP[SRC=10.10.20.2, DST=10.10.10.11] TCP[DPORT=30022]} hacia el router \texttt{r1} (\texttt{10.10.20.1}).
  \item \texttt{r1} encamina sin NAT entre \texttt{user\_LAN} y \texttt{k8s\_LAN}, entregando el paquete a \texttt{e1} de \texttt{k8s-worker1} con \texttt{SRC=10.10.20.2}.
  \item En \texttt{PREROUTING}, \texttt{kube-proxy} aplica DNAT local hacia \texttt{DST=PodAdminIP:22} \textbf{sin marcar \texttt{KUBE-MARK-MASQ}} (gracias a \texttt{Local}).
  \item En \texttt{FORWARD} hacia la interfaz \texttt{cali*} del Pod \texttt{admin}, Calico evalúa \texttt{allow-ssh-to-admin}: comprueba \texttt{SRC=10.10.20.2 $\in$ 10.10.20.2/32} y \texttt{DPORT=22/TCP} $\rightarrow$ \textbf{\texttt{ACCEPT}}. Cualquier intento desde otra máquina (como \texttt{k8s-master} \texttt{10.10.10.10} o un nodo no autorizado) llega con su IP real y recibe \textbf{\texttt{DROP}}.
\end{enumerate}

\section{Decisión 5: Aprovisionamiento de Herramientas bajo Bloqueo \texttt{Egress}}

\begin{decisionbox}{Problema: Instalación de \texttt{mariadb-client} cuando \texttt{default-deny-all} bloquea Internet}
Las imágenes \texttt{dgsgiros/admin-ssh} y \texttt{dgsgiros/nginx-ssh} incluyen utilidades de red y \texttt{sudo apt}, pero no traen preinstalado el cliente \texttt{mariadb-client} necesario en el Apartado 2.7 para administrar la base de datos desde los Pods. Una vez aplicada la política \texttt{default-deny-all} (que bloquea todo \texttt{Egress} en \texttt{eth0}), cualquier \texttt{apt-get update} directo hacia Internet por \texttt{eth0} queda bloqueado, incluso si un Pod web se recrea posteriormente al escalar el \texttt{Deployment}.
\end{decisionbox}

\textbf{Solución de ingeniería implementada:}
En el contenedor LXC \texttt{logs} (\texttt{10.1.2.101}) se despliega un proxy HTTP ligero (\textbf{\texttt{tinyproxy}} en el puerto \texttt{8888}, restringido mediante ACL \texttt{Allow 10.1.2.0/24}). De este modo, sin abrir ningún agujero en las \texttt{NetworkPolicies} de Kubernetes sobre \texttt{eth0}, tanto \texttt{01\_deploy\_all.sh} como el administrador pueden instalar paquetes en cualquier Pod canalizando \texttt{apt} por la interfaz secundaria \texttt{net1} (VLAN \texttt{1000}) hacia \texttt{10.1.2.101:8888}, saliendo a los repositorios oficiales de Ubuntu/Debian a través del enlace \texttt{eth9} (\texttt{virbr0}) del escenario VNX.

\section{Decisión 6: Diseño del Escenario VNX LXC (\texttt{k8s\_lxc\_vlan.xml})}

En \texttt{vnx/k8s\_lxc\_vlan.xml} se tuvieron en cuenta cuatro aspectos críticos de integración con VNX 2.0 y Linux:
\begin{enumerate}
  \item \textbf{\texttt{managed="no"} en \texttt{VlanNet} y \texttt{virbr0}}: impide que el ciclo de vida del escenario complementario intente recrear o destruir el switch Open vSwitch \texttt{VlanNet} que pertenece al escenario principal de Kubernetes.
  \item \textbf{Desplazamiento de MACs y red privada de gestión (\texttt{automac offset="10"}, \texttt{offset="30"})}: evita colisiones de direcciones MAC en Capa 2 y de direcciones \texttt{10.250.0.x} en el host con las máquinas virtuales del clúster K8s.
  \item \textbf{MariaDB escuchando en la VLAN (\texttt{bind-address = 0.0.0.0})}: en Ubuntu/Debian, el paquete \texttt{mariadb-server} escucha por defecto únicamente en \texttt{127.0.0.1}, rechazando con \texttt{Connection refused} cualquier cliente externo. Mediante \texttt{<filetree>} se inyecta \texttt{50-server.cnf} habilitando la escucha en red y \texttt{init\_db.sql} creando la base de datos \texttt{cnvr\_db} y el usuario administrador \texttt{'dbadmin'@'10.1.2.\%'} acotado criptográfica y topológicamente a la subred de la VLAN.
  \item \textbf{Recepción centralizada Rsyslog (\texttt{10-remote-logs.conf})}: activa los módulos \texttt{imudp} e \texttt{imtcp} en el puerto \texttt{514} del contenedor \texttt{logs}, clasificando automáticamente los mensajes entrantes por IP de origen en \texttt{/var/log/remotos/<IP>/<programa>.log}.
\end{enumerate}

\chapter{Verificación Experimental, Pruebas y Conclusiones}

\section{Flujo Automatizado de Despliegue (\texttt{01\_deploy\_all.sh})}

Todo el proceso de puesta en marcha se condensa en un único comando ejecutable desde el host del laboratorio:
\begin{lstlisting}[language=bash]
chmod +x scripts/*.sh
./scripts/01_deploy_all.sh
\end{lstlisting}

El script ejecuta de forma secuencial y limpia (\texttt{set -euo pipefail}) los seis pasos del escenario sobre el entorno del laboratorio:
\begin{enumerate}
  \item Generación de \texttt{keys/id\_rsa\_internal} (Par 1) y \texttt{keys/id\_rsa\_external} (Par 2) y copia de la clave pública interna hacia \texttt{vnx/conf/ssh/authorized\_keys}.
  \item Arranque del clúster Kubernetes con Calico (\path{/lab/cnvr/bin/get-k8s-calico} $\rightarrow$ \path{./bin/start-k8s-calico}), carga de credenciales con \path{./bin/get-master-kubeconf}, despliegue del nodo proveedor \textbf{\texttt{terraform}} (\texttt{sudo vnx -f} \path{k8s_lab-terraform.xml} \texttt{--create} y \texttt{-x install-terraform}) y exportación de su kubeconfig con \path{./bin/set-terraform-kubeconf}.
  \item Creación de la subinterfaz \texttt{e2.1000} y arranque del demonio CNI \texttt{dhcp} (\path{/run/cni/dhcp.sock}) en \texttt{k8s-worker1} y \texttt{k8s-worker2}.
  \item Despliegue del escenario complementario VNX (\texttt{sudo vnx -f k8s\_lxc\_vlan.xml -v --create} y \texttt{-x setup}) para aprovisionar OpenSSH, MariaDB, Rsyslog, Dnsmasq y Tinyproxy en \texttt{bbdd} y \texttt{logs}.
  \item Sincronización de \texttt{terraform/*.tf}, \texttt{terraform.tfvars} y \texttt{keys/} en \texttt{/shared/cnvr\_final/} y ejecución remota en el nodo \texttt{terraform} de \texttt{terraform init}, \texttt{terraform validate} y \texttt{terraform apply -auto-approve}, respaldando el fichero de estado.
  \item Instalación de \texttt{mariadb-client} en el Pod \texttt{admin} y en las réplicas web empleando el proxy interno de la VLAN (\texttt{10.1.2.101:8888}), mostrando finalmente \texttt{terraform output}.
\end{enumerate}

\section{Protocolo de Verificación Integral (\texttt{02\_verify\_all.sh})}

El script \texttt{scripts/02\_verify\_all.sh} automatiza las seis baterías de pruebas exigidas por el enunciado, verificando tanto el funcionamiento esperado (pruebas positivas) como el bloqueo efectivo de accesos prohibidos (pruebas negativas).

\subsection{Verificación de Interfaces Multus y Etiquetado 802.1Q en la VLAN}

Al inspeccionar las interfaces dentro del Pod \texttt{admin} y de las tres réplicas de \texttt{web-deployment} mediante \texttt{kubectl exec <pod> -- ip -4 addr}, se comprueba la presencia simultánea de:
\begin{itemize}
  \item \textbf{\texttt{eth0}}: dirección del pool CNI de Calico para el tráfico interno del clúster.
  \item \textbf{\texttt{net1}}: interfaz \texttt{macvlan} en la subred \texttt{10.1.2.0/24} (\texttt{10.1.2.10/24} estática en \texttt{admin}, y concesiones DHCP \texttt{10.1.2.50/24}, \texttt{.51/24}, \texttt{.52/24} en los Pods web).
\end{itemize}

Al capturar tráfico con \texttt{tcpdump -i e2 -n -e icmp} en \texttt{k8s-worker1} durante los pings hacia \texttt{10.1.2.100} (\texttt{bbdd}) y \texttt{10.1.2.101} (\texttt{logs}), se observa que las tramas abandonan el nodo \textit{worker} directamente con la dirección MAC propia de la subinterfaz \texttt{macvlan} del Pod y encapsuladas con cabecera \textbf{IEEE 802.1Q (\texttt{ethertype 0x8100, vlan 1000})}, sin sobrecarga de cabeceras VXLAN ni IPIP, y con \texttt{TTL=64} tanto en la ida como en la vuelta al encontrarse todos los elementos en el mismo dominio de difusión L2.

\subsection{Verificación de Acceso SSH Cruzado mediante Pares de Claves}

\begin{enumerate}
  \item \textbf{Desde el Host (\texttt{user\_LAN}) al Pod \texttt{admin} (Par 2)}:
\begin{lstlisting}[language=bash]
ssh -i keys/id_rsa_external -p 30022 adminuser@10.10.10.11
\end{lstlisting}
  El acceso se completa sin solicitar contraseña a través del \texttt{NodePort} de \texttt{k8s-worker1}, autenticando contra el volumen montado desde el \texttt{Secret} \texttt{admin-ssh-authkeys}.
  \item \textbf{Desde el Pod \texttt{admin} hacia los Pods Web y Contenedores LXC (Par 1)}:
  Una vez dentro de \texttt{admin} como \texttt{adminuser}, las conexiones \texttt{ssh} hacia \path{webuser@10.1.2.5X}, \path{adminuser@10.1.2.100} (\texttt{bbdd}) y \path{adminuser@10.1.2.101} (\texttt{logs}) se establecen sin contraseña empleando la clave privada \texttt{\textasciitilde/.ssh/id\_rsa} inyectada por \texttt{admin-ssh-privkey}. Asimismo, cualquier intento de inicio de sesión por contraseña es rechazado por \texttt{sshd} (\texttt{PasswordAuthentication no}).
\end{enumerate}

\subsection{Administración de Base de Datos MariaDB y Recepción de Logs}

Desde el Pod \texttt{admin} y desde cualquiera de los Pods de servidores web se verifica la administración completa (lectura, inserción de registros y gestión de esquemas) sobre el servidor MariaDB en \texttt{bbdd} (\texttt{10.1.2.100:3306}):
\begin{lstlisting}[language=bash]
mariadb -h 10.1.2.100 -u dbadmin -pcnvr2026 cnvr_db \
  -e "INSERT INTO nodos_registrados (nombre, rol, ip_vlan)
      VALUES ('$(hostname)', 'Pod Web Nginx', '10.1.2.50');
      SELECT * FROM nodos_registrados;"
\end{lstlisting}
De forma análoga, el envío de eventos syslog hacia \texttt{10.1.2.101:514} registra instantáneamente las entradas clasificadas por IP de Pod en \texttt{/var/log/remotos/10.1.2.10/} dentro del contenedor LXC \texttt{logs}.

\subsection{Personalización Web, Balanceo Ingress, Auto-reparación y Escalado}

Al realizar peticiones sucesivas desde el host al controlador Ingress en la ruta del Grupo 11:
\begin{lstlisting}[language=bash]
for i in {1..6}; do
  curl -s http://10.10.10.20/servers-G11 | grep "Pod (Servidor)"
done
\end{lstlisting}
Las respuestas alternan entre las tres réplicas del \texttt{Deployment}, mostrando cada una su identificador único sustituido desde la plantilla del \texttt{ConfigMap} por su respectivo \texttt{initContainer}. Si se elimina manualmente un Pod (\texttt{kubectl delete pod <web-pod>}), el \texttt{ReplicaSet} instancia inmediatamente un sustituto con nuevo sufijo, nueva IP DHCP en la VLAN \texttt{1000} y su página HTML recién renderizada; y si se ejecuta \texttt{WEB\_REPLICAS=5 ./scripts/01\_deploy\_all.sh} (o \texttt{terraform apply -var="web\_replicas=5"} desde \texttt{/shared/cnvr\_final/terraform} en el nodo \texttt{terraform}), el servicio escala limpiamente a 5 instancias balanceadas y actualiza el output \texttt{web\_pod\_names} sin pérdida de peticiones.

\subsection{Matriz de Auditoría de \texttt{NetworkPolicies} (Calico)}

La Tabla~\ref{tab:matriz-seguridad} resume las pruebas positivas y negativas verificadas por \texttt{02\_verify\_all.sh} tras aplicar las tres políticas de red:

\begin{table}[H]
\centering
\footnotesize
\renewcommand{\arraystretch}{1.25}
\begin{tabularx}{\textwidth}{l l l c >{\raggedright\arraybackslash}X}
\toprule
\textbf{Origen} & \textbf{Destino} & \textbf{Plano / Puerto} & \textbf{Resultado} & \textbf{Política Responsable} \\
\midrule
Host (\texttt{10.10.20.2}) & Ingress $\rightarrow$ Web & HTTP \texttt{:80} (\texttt{eth0}) & \textcolor{okgreen}{\textbf{PERMITIDO}} & \path{allow-ingress-to-web} (AND \texttt{ns} + \texttt{pod}) \\
Pod \texttt{admin} / otro Pod & IP Calico Pod Web & HTTP \texttt{:80} (\texttt{eth0}) & \textcolor{errred}{\textbf{BLOQUEADO}} & \path{default-deny-all} (no es \texttt{ingress-nginx}) \\
Host (\texttt{10.10.20.2}) & \texttt{10.10.10.11:30022} & SSH \texttt{:22} (\texttt{eth0}) & \textcolor{okgreen}{\textbf{PERMITIDO}} & \path{allow-ssh-to-admin} (\texttt{ipBlock} + \texttt{Local}) \\
\texttt{k8s-master} (\texttt{.10.10}) & \texttt{10.10.10.11:30022} & SSH \texttt{:22} (\texttt{eth0}) & \textcolor{errred}{\textbf{BLOQUEADO}} & Fuera de \texttt{10.10.20.2/32} (IP preservada) \\
Pods \texttt{admin} y \texttt{web} & Internet (\texttt{1.1.1.1}) & Salida (\texttt{eth0}) & \textcolor{errred}{\textbf{BLOQUEADO}} & \path{default-deny-all} (bloqueo \texttt{Egress}) \\
Pods \texttt{admin} y \texttt{web} & Web / \texttt{bbdd} / \texttt{logs} & VLAN \texttt{1000} (\texttt{net1}) & \textcolor{okgreen}{\textbf{PERMITIDO}} & Plano L2 \texttt{macvlan} aislado en \texttt{VlanNet} \\
\bottomrule
\end{tabularx}
\caption{Matriz de resultados de seguridad y conectividad verificados en el escenario.}
\label{tab:matriz-seguridad}
\end{table}

\section{Conclusiones}

El desarrollo de este trabajo final demuestra cómo integrar de forma limpia y totalmente automatizada el mundo de las infraestructuras tradicionales/virtualizadas en Linux (\textbf{VNX}, contenedores de sistema \textbf{LXC}, conmutación \textbf{Open vSwitch} y etiquetado \textbf{VLAN 802.1Q}) con arquitecturas modernas de microservicios sobre \textbf{Kubernetes} (\textbf{Calico}, \textbf{Multus}, \textbf{MetalLB} e \textbf{Ingress NGINX}) bajo el gobierno declarativo de \textbf{Terraform}.

Más allá del cumplimiento estricto de todos los requisitos funcionales del enunciado, resolver problemas sutiles de sistemas ---como la coordinación IPAM mediante DHCP sin puerta de enlace en Multus para permitir escalabilidad real en \texttt{Deployments}, la conciliación entre \texttt{fsGroup: 1000} de Kubernetes y \texttt{StrictModes} de OpenSSH mediante \texttt{initContainers} sin privilegios, y la preservación de la IP origen del cliente mediante \texttt{externalTrafficPolicy: Local} para habilitar el filtrado \texttt{ipBlock} en Calico--- pone de manifiesto que dominar las abstracciones de alto nivel de la nube pasa inevitablemente por comprender con precisión cómo operan los \textit{namespaces}, interfaces virtuales, túneles IP-in-IP y tablas de \texttt{netfilter} del kernel de Linux que las sustentan.

\begin{thebibliography}{10}
\bibitem{tf-k8s} HashiCorp, \textit{Terraform Kubernetes Provider Documentation}, \url{https://registry.terraform.io/providers/hashicorp/kubernetes/latest/docs}.
\bibitem{tf-manifest} HashiCorp, \textit{Resource: kubernetes\_manifest --- Server-Side Apply \& CRDs}, \url{https://registry.terraform.io/providers/hashicorp/kubernetes/latest/docs/resources/manifest}.
\bibitem{tf-os} OpenStack / HashiCorp, \textit{Terraform OpenStack Provider Documentation}, \url{https://registry.terraform.io/providers/terraform-provider-openstack/openstack/latest/docs}.
\bibitem{k8s-secrets} Kubernetes Documentation, \textit{Secrets \& Distribute Credentials Securely Using Secrets}, \url{https://kubernetes.io/docs/concepts/configuration/secret/}.
\bibitem{k8s-netpol} Kubernetes Documentation, \textit{Network Policies \& Source IP Preservation for Services with Type=NodePort}, \url{https://kubernetes.io/docs/concepts/services-networking/network-policies/}.
\bibitem{calico} Tigera / Project Calico, \textit{Calico Networking (IP-in-IP, BGP) and Network Policy Architecture}, \url{https://docs.tigera.io/calico/latest/about/}.
\bibitem{multus} K8s Network Plumbing Working Group, \textit{Multus-CNI: Multi-homed pods in Kubernetes}, \url{https://github.com/k8snetworkplumbingwg/multus-cni}.
\bibitem{vnx} D.~Fernández (DIT-UPM), \textit{Virtual Networks over LinuX (VNX) Documentation \& XML 2.0 Specification}, \url{http://web.dit.upm.es/vnxwiki/}.
\bibitem{brikman} Y.~Brikman, \textit{Terraform: Up and Running: Writing Infrastructure as Code}, 3rd~ed., O'Reilly Media, 2022.
\end{thebibliography}


\end{document}
